\documentclass[11pt]{article}

\usepackage[a4paper,margin=27mm]{geometry}
\usepackage{amsmath,amssymb,amsthm,mathtools}
\usepackage{array}
\usepackage{booktabs}
\usepackage{graphicx}
\graphicspath{{../conference-paper-2/}}
\usepackage[T1]{fontenc}
\usepackage{lmodern}
\usepackage{microtype}
\usepackage{xcolor}
\usepackage[hidelinks]{hyperref}

\newcommand{\C}{\mathbb C}
\newcommand{\Z}{\mathbb Z}
\newcommand{\T}{\mathbb T}
\newcommand{\norm}[1]{\left\lVert #1\right\rVert}
\newtheorem{theorem}{Theorem}[section]
\newtheorem{proposition}[theorem]{Proposition}
\newtheorem{corollary}[theorem]{Corollary}

\title{\textbf{Spectral Equivalence, Hamiltonian Fidelity, and Hopping Locality\\
in Two- and Three-Dimensional Rank-Two Models}}
\author{Eugene Joseph M. Ragasa\\
Department of Physics, De La Salle University, Manila, Philippines}
\date{}

\begin{document}
\maketitle

\begin{abstract}
Band eigenvalues do not determine a matrix-valued tight-binding representation or its
real-space locality. We formulate a rank-two framework that keeps the finite
cell--orbital, Bloch, spectral, retained-frame, and Wannier representations explicit and
aligns their state spaces before comparing Hamiltonians. A Laurent-polynomial
obstruction proves that spectral projection does not preserve finite hopping support in
general. For orthogonal flat-parent models, exact Bessel identities equate the
irreducible restricted-alignment residual with the normalized Fourier weight outside the
onsite block, while shell-resolved Jacobi--Anger formulas provide factorial tail bounds.
Bond-orbit identities extend the analytic comparison to non-orthogonal but physically
isotropic lattices using complete metric shells. Calculations for eight square,
triangular, honeycomb, simple-cubic, FCC, and diamond-inspired cases reveal non-collapsed
sampled Pareto sets in retained-block count, Hamiltonian error, and sorted-band error.
The optimized comparison scale ranges from $0.729890$ to $0.937745$; at the largest
diamond shell, restoring unit scale increases the Hamiltonian mismatch by about $2.9\%$
while reducing the sorted-band mismatch by about $82\%$. Synthetic rhombic and Bain
paths separate geometry-only reordering, a frozen-parent null response, and a prescribed
radial response. Across 201 samples, the retained subspaces remain isolated and
well-conditioned, the frozen-parent response stays at binary64 roundoff, and the radial
response is traceless dominated. Exact gauge equivalence therefore coexists with a
case-dependent three-way compromise once one restricts alignment and hopping support;
the present comparison does not isolate metric-shell geometry from parent connectivity.
The numerical statements are replaceable finite-mesh development results for synthetic
models, not qualified retained evidence, continuum certificates, or material predictions.
\end{abstract}

\section{Introduction}

Pointwise band equivalence does not determine hopping locality under a fixed
identification of retained coordinates. Here we extend that distinction across spatial
dimension and lattice geometry while the retained rank remains two. The purpose is not
to infer a Hamiltonian from eigenvalues alone. It is to compare
complete Hamiltonian representations connected by declared isometries and gauges, and
to quantify the representation cost introduced by finite hopping support.

The gauge dependence and localization of Wannier representations are established
subjects, from analytic and exponential localization results
\cite{Kohn1959AnalyticWannier,Brouder2007ExponentialWannier} through practical
maximally localized constructions and Wannier90
\cite{Mostofi2008Wannier90,Marzari2012MaximallyLocalized}. Compact support is a
stronger algebraic property and need not coincide with maximal localization
\cite{Sathe2024CompactWannier}. Material-specific Wannier and strained-Hamiltonian
models likewise exist for graphene and other hexagonal systems
\cite{Jung2013GrapheneWannier,Fang2018Strained2D}, while environment-dependent
Hamiltonian mappings address geometry-responsive electronic structure
\cite{Zhang2022EquivariantHamiltonians}. Our contribution is the controlled combination
of pointwise spectral equivalence, declared gauge alignment, metric-complete truncation,
and a decision landscape relating hopping complexity to operator and spectral error.

The study addresses three questions:
\begin{enumerate}
 \item How are a finite cell--orbital operator, its Bloch representation, its complete
       eigensystem, a retained periodic frame, and its Wannier/tight-binding blocks
       related without conflating their state spaces?
 \item After those spaces have been aligned, what compromises among retained hopping
       blocks, Hamiltonian agreement, and sorted-band agreement remain when alignment is
       restricted to one cell-local rotation and hopping support is finite?
 \item Which parts of that compromise arise from physical-metric shell reordering,
       which vanish for an unchanged parent, and which respond to a prescribed radial
       change of bond amplitudes along synthetic deformation paths?
\end{enumerate}

The analysis yields two theorem-level results and two finite sampled development
findings. First, finite hopping support is not preserved by spectral projection in
general. Second, the orthogonal flat-parent family admits exact residual--locality
identities and certified Fourier tails. In the replaceable development calculations,
all eight fixed-geometry cases exhibit a sampled trade-off rather than a single optimum
shared by locality, operator fidelity, and spectral fidelity. The sampled strain
controls separate metric reordering from Hamiltonian response: the frozen parent is a
numerical null test, whereas the radial response is finite, predominantly traceless, and
unaccompanied by loss of retained-subspace isolation or frame conditioning. These
numerical findings are not qualified retained evidence.

The scope is deliberately rank two. Square and simple-cubic systems establish
orthogonal baselines; triangular/honeycomb and FCC/diamond-inspired systems test
non-orthogonal but physically isotropic geometry. Direction-dependent physical coupling
laws, higher-rank $U(r)$ gauges with $r>2$, and material parametrizations lie outside the
present study.

\section{Theoretical framework and methods}

This section fixes the spaces and maps used throughout the paper, following the finite
spectral and tight-binding realization defined in the orthogonal family protocols.

\subsection{Finite periodic spaces and transforms}

Let $d\in\{2,3\}$ be the spatial dimension and let $N$ be an even number of periodically
identified cells along each primitive direction. The finite cell set and its centered
reduced reciprocal mesh are
\begin{align}
 \Lambda_N^{(d)}&=(\mathbb Z/N\mathbb Z)^d,\\
 \mathcal Q_N^{(d)}
 &=\left\{\mathbf q_{\mathbf j}:q_{j_\alpha}
   =\frac{j_\alpha-N/2}{N},\quad
   j_\alpha\in\{0,\ldots,N-1\}\right\}.
 \label{eq:finite-cell-and-mesh}
\end{align}
A vector $\mathbf n\in\Lambda_N^{(d)}$ labels a periodic cell, whereas
$\mathbf q\in\mathcal Q_N^{(d)}\subset[-1/2,1/2)^d$ is a dimensionless reduced
reciprocal coordinate. The Bloch phase pairing is
$e^{2\pi i\mathbf q\cdot\mathbf n}$.

The cell-amplitude space and the two operator state spaces are
\begin{equation}
 \mathcal L_N^{(d)}
 =\operatorname{span}\{\lvert\mathbf n\rangle:\mathbf n\in\Lambda_N^{(d)}\}
 \cong(\C^N)^{\otimes d},\qquad
 \mathcal H_{p,N}^{(d)}=\mathcal L_N^{(d)}\otimes\C^4,
 \qquad
 \mathcal H_{r,N}^{(d)}=\mathcal L_N^{(d)}\otimes\C^2.
 \label{eq:spaces}
\end{equation}
The subscript $p$ denotes the four-state parent orbital fiber and $r$ denotes the
rank-two retained fiber. Thus $\widehat K_{\rho,N}^{(d)}$ acts on
$\mathcal H_{p,N}^{(d)}$, while the retained Bloch and Wannier operators act on
$\mathcal H_{r,N}^{(d)}$. The symbol $I_m$ denotes the identity on $\C^m$.

Let $\mathsf F_N$ be the unitary centered discrete Fourier transform and define the
one-axis cyclic translation $\mathsf S_N$ by
\begin{equation}
 \mathsf F_N\mathsf S_N\mathsf F_N^\dagger
 =\operatorname{diag}_{q\in\mathcal Q_N^{(1)}}e^{2\pi iq}.
 \label{eq:one-axis-translation}
\end{equation}
The translation on primitive direction $\alpha$ and the $d$-dimensional Fourier
transform are
\begin{equation}
 \mathsf S_{\alpha,N}
 =I_N^{\otimes(\alpha-1)}\otimes\mathsf S_N
  \otimes I_N^{\otimes(d-\alpha)},
 \qquad
 \mathcal F_{d,N}=\mathsf F_N^{\otimes d}.
 \label{eq:tensor-translation-fourier}
\end{equation}
These Kronecker products act only on cell labels. They do not tensor the
four-dimensional orbital fiber.

\subsection{Exact orthogonal parent}

Let $E_{\mathrm{ref}}>0$ be the fixed energy unit; it is not a band gap. The orthogonal
directed-parent cases use the protocol matrices
\begin{equation}
 T_0=E_{\mathrm{ref}}
 \begin{pmatrix}
 -1.20&0.06&0&0\\
 0.06&-0.35&0&0\\
 0&0&1.00&0.04\\
 0&0&0.04&1.75
 \end{pmatrix},
 \label{eq:exact-onsite-block}
\end{equation}
\begin{equation}
 T_+=E_{\mathrm{ref}}
 \begin{pmatrix}
 0.15&0&0.12&0\\
 0&-0.10&0&0.10\\
 0&0&0.05&0\\
 0&0&0&-0.04
 \end{pmatrix},
 \qquad T_-=T_+^\dagger.
 \label{eq:exact-directed-hopping-block}
\end{equation}
Thus $T_0\in\C^{4\times4}$ is the Hermitian onsite block,
$T_+\in\C^{4\times4}$ is the directed positive-axis hopping block, and $T_-$ is its
reverse. The parameter $\rho\in[0,1]$ is a dimensionless parent-hopping homotopy: at
$\rho=0$ the parent is cell local, while $\rho=1$ gives the declared dispersive parent.
The finite cell--orbital operator is
\begin{equation}
 \widehat K_{\rho,N}^{(d)}
 =I_{N^d}\otimes T_0+
 \frac{\rho}{d}\sum_{\alpha=1}^d
 \left(\mathsf S_{\alpha,N}\otimes T_+
 +\mathsf S_{\alpha,N}^\dagger\otimes T_-\right).
 \label{eq:kronecker-parent}
\end{equation}
The factor $1/d$ keeps the summed axis-hopping scale fixed when dimension changes.
Fourier transformation of the cell coordinate gives the direct sum
\begin{align}
 \widehat K_{B,\rho,N}^{(d)}
 &:=(\mathcal F_{d,N}\otimes I_4)\widehat K_{\rho,N}^{(d)}
 (\mathcal F_{d,N}^\dagger\otimes I_4)\notag\\
 &=\bigoplus_{\mathbf q\in\mathcal Q_N^{(d)}}
 K_\rho^{(d)}(\mathbf q),
 \label{eq:bloch-map}\\
 K_\rho^{(d)}(\mathbf q)
 &=T_0+\frac{\rho}{d}\sum_{\alpha=1}^d
 \left(e^{2\pi iq_\alpha}T_+
       +e^{-2\pi iq_\alpha}T_-\right).
 \label{eq:bloch-fiber-definition}
\end{align}
Here $B$ denotes the Bloch representation, and each
$K_\rho^{(d)}(\mathbf q)\in\C^{4\times4}$ is the parent Hamiltonian at one reciprocal
mesh point.

\subsection{Complete spectrum and retained frame}

At each $\mathbf q$, a complete eigensystem is
\begin{equation}
 K_\rho^{(d)}(\mathbf q)
 =Q(\mathbf q)E(\mathbf q)Q^\dagger(\mathbf q),
 \qquad
 E(\mathbf q)=\operatorname{diag}
 (\varepsilon_1(\mathbf q),\ldots,\varepsilon_4(\mathbf q)),
 \label{eq:eigensystem}
\end{equation}
where $Q(\mathbf q)\in U(4)$ contains the corresponding eigenvectors and the real
eigenvalues are ordered increasingly. The diagonal array $E(\mathbf q)$ is an operator
only in its eigenvector coordinates. Without $Q(\mathbf q)$ it is not an operator in the
parent orbital basis and cannot be subtracted from a tight-binding Hamiltonian.

Let $I=\{1,2\}$ denote the retained two lowest bands. Their spectral projector is
\begin{equation}
 P_I(\mathbf q)=Q_I(\mathbf q)Q_I^\dagger(\mathbf q),
 \label{eq:retained-projector}
\end{equation}
where $Q_I$ contains the two selected eigenvectors. A periodic retained frame is an
isometry $F(\mathbf q):\C^2\to\C^4$ that satisfies
\begin{equation}
 F^\dagger(\mathbf q)F(\mathbf q)=I_2,
 \qquad
 F(\mathbf q)F^\dagger(\mathbf q)=P_I(\mathbf q),
 \label{eq:frame-definition}
\end{equation}
with the declared reciprocal periodicity and transport qualification. The retained
projector is admissible only while bands two and three remain externally isolated; a
continuous lower bound is distinguished from a sampled minimum.

The frame construction freezes an anchor, mesh traversal, neighbor-to-neighbor unitary
Procrustes rule, reciprocal-boundary sewing rule, path-consistency diagnostics, and rank
thresholds. If a periodic full-rank frame cannot be constructed and verified, the
protocol stops rather than selecting a gauge across an obstruction. Transport closure
and plaquette phases are sewing diagnostics here, not topological claims. Projectors and
reconstructed operators provide cross-mesh evidence; raw frame entries do not.

The frame chooses coordinates inside the retained subspace. Changing it by a
momentum-dependent $U(2)$ rotation changes the hopping representation while preserving
the retained bands. The retained Bloch matrix is
\begin{equation}
 H(\mathbf q)=F^\dagger(\mathbf q)
 K_\rho^{(d)}(\mathbf q)F(\mathbf q)\in\C^{2\times2}.
 \label{eq:retained}
\end{equation}
\subsection{Retained Bloch and Wannier operators}

With $\mathcal V=\bigoplus_{\mathbf q}F(\mathbf q)$, the block isometry maps the retained
Bloch direct sum into the parent-orbital Bloch direct sum. The complete retained Bloch
operator and its cell-localized Wannier representation are
\begin{align}
 \widehat H_B
 &=\mathcal V^\dagger\widehat K_{B,\rho,N}^{(d)}\mathcal V
 =\bigoplus_{\mathbf q}H(\mathbf q),\\
 \widehat H_W
 &=(\mathcal F_{d,N}^\dagger\otimes I_2)\widehat H_B
   (\mathcal F_{d,N}\otimes I_2).
 \label{eq:retained-wannier-operator}
\end{align}
The $2\times2$ block connecting two cells separated by cyclic mode
$\mathbf n\in\Lambda_N^{(d)}$ is the finite Fourier hopping block
\begin{equation}
 T_{\mathbf n}^{(N)}
 =\frac{1}{N^d}\sum_{\mathbf q\in\mathcal Q_N^{(d)}}
 e^{-2\pi i\mathbf q\cdot\mathbf n}H(\mathbf q),
 \qquad
 H(\mathbf q)=\sum_{\mathbf n\in\Lambda_N^{(d)}}
 e^{2\pi i\mathbf q\cdot\mathbf n}T_{\mathbf n}^{(N)}.
 \label{eq:finite-hopping-transform}
\end{equation}
``Complete'' means that every cyclic Fourier mode is retained. Interpolation onto a new
mesh, conversion to an ordinary conjugate-closed inventory, and deliberate
physical-radius truncation occur only after this exact finite representation chain.

\subsection{Spectral projection is not finite-range compression}

For a spectral projector $P_I$ with isometry $V_I$ onto its range,
\begin{equation}
 H_{\mathrm{red}}=V_I^\dagger H V_I,
 \qquad
 V_IH_{\mathrm{red}}V_I^\dagger=P_IHP_I.
 \label{eq:spectral-projection}
\end{equation}
This operation preserves the retained eigenpairs exactly, but it need not preserve a
finite hopping range. The following elementary example makes that statement exact.

\begin{proposition}[Spectral projection can destroy finite hopping support]
\label{prop:projection-destroys-range}
Let
\begin{equation}
 H(k)=\begin{pmatrix}
 \Delta&t(1+e^{-ik})\\
 t(1+e^{ik})&-\Delta
 \end{pmatrix},
 \qquad \Delta t\ne0.
 \label{eq:projection-example}
\end{equation}
Then $H$ has nearest-neighbor Fourier support and eigenvalues
$E_\pm(k)=\pm d(k)$, where
\begin{equation}
 d(k)=\sqrt{\Delta^2+4t^2\cos^2(k/2)}.
\end{equation}
The embedded lower-band operator
\begin{equation}
 P_-(k)H(k)P_-(k)=\tfrac12[H(k)-d(k)I]
 \label{eq:projected-range}
\end{equation}
is not a matrix-valued finite Fourier polynomial.
\end{proposition}

\noindent\emph{Proof.} See Appendix~\ref{app:proof-projection-range}.


The gap condition $\Delta\ne0$ keeps the two bands separated. The proposition does not
claim that every spectral projection of every finite-range parent is long ranged; it
shows that finite support is not preserved as an operation. Projection, frame selection,
Fourier conversion, hopping truncation, and restriction to a candidate model class are
therefore distinct operations.

\section{Orthogonal isotropic baselines}

The square and simple-cubic baselines use
\begin{equation}
 K_\rho^{(d)}(\mathbf q)=T_0+\frac{\rho}{d}
 \sum_{\alpha=1}^d
 \left(e^{2\pi iq_\alpha}T_+
 +e^{-2\pi iq_\alpha}T_-\right)
 \label{eq:orthogonal-parent}
\end{equation}
and
\begin{equation}
 U_a^{(d)}(\mathbf q)=R\!\left[
 \frac{a}{d}\sum_{\alpha=1}^d\sin(2\pi q_\alpha)
 \right].
 \label{eq:rotation}
\end{equation}
The flat-parent correlations are $J_0(a)^2$ in two dimensions and
$J_0(2a/3)^3$ in three dimensions. The next theorem proves these formulas and their
exact real-space meaning. They are continuum analytic comparators and do not replace
finite-mesh correlations. Square and cubic truncations retain every attainable metric
shell through the declared cap using complete signed-permutation orbits.

\subsection{Flat-parent residual--locality theorem}

Throughout, $E_{\mathrm{ref}}>0$ is an arbitrary reference-energy normalization, not a
band gap. Give $\T^d=[-1/2,1/2)^d$ normalized Lebesgue measure. Let
$D\in\mathbb R^{2\times2}$ be nonzero, symmetric, and traceless, and let the flat
retained parent be $H(\mathbf q)=\bar E I_2+D$. For the rotation in
Eq.~\eqref{eq:rotation}, define
\begin{align}
 B_a^{(d)}(\mathbf q)
 &=U_a^{(d)}(\mathbf q)^\dagger D U_a^{(d)}(\mathbf q),\\
 B_{\mathbf R}^{(d)}
 &=\int_{\T^d}e^{-2\pi i\mathbf q\cdot\mathbf R}
 B_a^{(d)}(\mathbf q)\,d\mathbf q,
 \qquad \mathbf R\in\Z^d,\\
 \kappa&=\frac{\norm{D}_F^2}{2E_{\mathrm{ref}}^2}.
 \label{eq:flat-definitions}
\end{align}
For one momentum-independent real rotation $R(\phi)$, centroid shift $s$, and the
nonnegative closure $\lambda\ge0$ of the splitting-scale domain, define
\begin{equation}
 L_{\mathrm{op},\phi}^2(s,\lambda)
 =\frac{1}{2E_{\mathrm{ref}}^2}\int_{\T^d}
 \norm{sE_{\mathrm{ref}}I_2+\lambda R(\phi)^\dagger D R(\phi)
       -B_a^{(d)}(\mathbf q)}_F^2\,d\mathbf q.
 \label{eq:flat-operator-loss}
\end{equation}
The normalized best global-rotation correlation is
\begin{equation}
 \mu_{\max}^{(d)}(a)=
 \max_{\phi\in\mathbb R}
 \frac{1}{\norm D_F^2}
 \int_{\T^d}
 \operatorname{tr}\!\left[
   R(\phi)^\dagger D R(\phi)B_a^{(d)}(\mathbf q)
 \right]d\mathbf q.
 \label{eq:flat-mu-definition}
\end{equation}

\begin{theorem}[Multidimensional flat-parent residual--locality identity]
\label{thm:flat-parent}
For $d\in\{2,3\}$ and $0\le a\le\pi/4$, set
\begin{equation}
 c_d(a)=J_0\!\left(\frac{2a}{d}\right)^d.
 \label{eq:flat-correlation-general}
\end{equation}
Then $c_d(a)>0$ and
\begin{align}
 \int_{\T^d}e^{2i\theta_a^{(d)}(\mathbf q)}\,d\mathbf q
 &=c_d(a),
 \label{eq:bessel-factorization}\\
 B_{\mathbf0}^{(d)}&=c_d(a)D.
 \label{eq:onsite-average}
\end{align}
The best global-rotation correlation is
$\mu_{\max}^{(d)}(a)=c_d(a)$, and the exact restricted-alignment residual satisfies
\begin{align}
 \inf_{s\in\mathbb R,\,\lambda\ge0,\,\phi\in\mathbb R}
 L_{\mathrm{op},\phi}^2(s,\lambda)
 &=\kappa\left[1-c_d(a)^2\right]
 \label{eq:flat-residual}\\
 &=\frac{1}{2E_{\mathrm{ref}}^2}
   \sum_{\mathbf R\ne\mathbf0}\norm{B_{\mathbf R}^{(d)}}_F^2.
 \label{eq:flat-tail}
\end{align}
If exact spectral agreement additionally fixes $s=0$ and $\lambda=1$, then
\begin{equation}
 \inf_{\phi\in\mathbb R}L_{\mathrm{op},\phi}^2(0,1)
 =2\kappa\left[1-c_d(a)\right].
 \label{eq:flat-exact-spectrum}
\end{equation}
Thus the two-dimensional formulas are
$c_2(a)=J_0(a)^2$ and the three-dimensional formulas are
$c_3(a)=J_0(2a/3)^3$.
\end{theorem}

\noindent\emph{Proof.} See Appendix~\ref{app:proof-flat-parent}.


\begin{theorem}[Shell-resolved coefficients and factorial Fourier-tail bound]
\label{thm:orthogonal-fourier-tail}
Use a constant real rotation that diagonalizes $D$ as
$D=\delta\sigma_z$, and set
\begin{equation}
 z=\frac{2a}{d},\qquad
 j_{\mathbf R}(z)=\prod_{\alpha=1}^d J_{R_\alpha}(z),\qquad
 p(\mathbf R)=\sum_{\alpha=1}^dR_\alpha\pmod 2.
\end{equation}
Then every Fourier block is known exactly:
\begin{equation}
 B_{\mathbf R}^{(d)}
 =\delta j_{\mathbf R}(z)
 \begin{cases}
  \sigma_z,&p(\mathbf R)=0,\\
  i\sigma_x,&p(\mathbf R)=1.
 \end{cases}
 \label{eq:orthogonal-fourier-block}
\end{equation}
For the complete squared-radius shell $s\in\mathbb N_0$, define
\begin{equation}
 W_{d,s}(z)=
 \sum_{\substack{\mathbf R\in\mathbb Z^d\\\norm{\mathbf R}_2^2=s}}
 \prod_{\alpha=1}^dJ_{R_\alpha}(z)^2.
 \label{eq:shell-weight}
\end{equation}
Its normalized hopping-block power and the total shell inventory satisfy
\begin{align}
 \frac{1}{2E_{\mathrm{ref}}^2}
 \sum_{\norm{\mathbf R}_2^2=s}\norm{B_{\mathbf R}^{(d)}}_F^2
 &=\kappa W_{d,s}(z),
 \label{eq:shell-power}\\
 \sum_{s=0}^{\infty}W_{d,s}(z)&=1.
 \label{eq:shell-partition}
\end{align}
For an integer shell cap $S\ge0$, let
\begin{equation}
 L=\left\lfloor\sqrt{\frac{S}{d}}\right\rfloor+1,
 \qquad x=\frac{|z|}{2}.
\end{equation}
Whenever $x<L+1$, the omitted normalized power has the rigorous bound
\begin{equation}
 \frac{1}{2E_{\mathrm{ref}}^2}
 \sum_{\norm{\mathbf R}_2^2>S}\norm{B_{\mathbf R}^{(d)}}_F^2
 \le
 \kappa\,
 \frac{2d\,e^{z^2/2}x^{2L}}
 {(L!)^2\left[1-x^2/(L+1)^2\right]}.
 \label{eq:orthogonal-factorial-tail}
\end{equation}
Thus the orthogonal flat-parent Fourier tail is certified directly from $a$, $d$, and
$S$; no fitted decay law is required.
\end{theorem}

\noindent\emph{Proof.} See Appendix~\ref{app:proof-orthogonal-tail}.


\begin{corollary}[Monotonicity and the small-amplitude dimension law]
\label{cor:dimension-law}
Define the minimum-alignment and exact-spectrum thresholds by
\begin{equation}
 \tau_{\min,d}^2(a)=\kappa[1-c_d(a)^2],
 \qquad
 \tau_{\mathrm{exact},d}^2(a)=2\kappa[1-c_d(a)].
\end{equation}
Both thresholds are strictly increasing for $0<a\le\pi/4$. As $a\to0$,
\begin{align}
 c_d(a)
 &=1-\frac{a^2}{d}+\frac{2d-1}{4d^3}a^4+O(a^6),
 \label{eq:correlation-small-a}\\
 \tau_{\min,d}^2(a)
 &=\frac{2\kappa}{d}a^2
   -\frac{\kappa(4d-1)}{2d^3}a^4+O(a^6),
 \label{eq:min-threshold-small-a}\\
 \tau_{\mathrm{exact},d}^2(a)
 &=\frac{2\kappa}{d}a^2
   -\frac{\kappa(2d-1)}{2d^3}a^4+O(a^6).
 \label{eq:exact-threshold-small-a}
\end{align}
In particular, both thresholds have leading slope
$\sqrt{2\kappa/d}$ with respect to the maximum rotation amplitude $a$.
\end{corollary}

\noindent\emph{Proof.} See Appendix~\ref{app:proof-dimension-law}.


Theorems~\ref{thm:flat-parent} and \ref{thm:orthogonal-fourier-tail}, together with
Corollary~\ref{cor:dimension-law}, apply to the orthogonal flat-parent master family.
The factorial bound certifies its complete-shell Fourier tail, but it must not be
transferred to a dependent non-orthogonal bond orbit. The bond-orbit theorem below
supplies that comparator's analytic form; a lattice-specific constrained-series tail
still requires a separate certificate.

The two- and three-dimensional calculations use independent meshes and shell
inventories. Their numerical outcomes are compared only through the common normalized
diagnostics defined above; no shell index or convergence statement is transferred
between dimensions.

\subsection{Two-dimensional scalar precursor}

An earlier retained scalar calculation used
\begin{equation}
 \widehat H_{2\mathrm D}=-\nabla^2+0.5\cos x+0.5\cos y
 +\lambda_{xy}\cos x\cos y.
 \label{eq:legacy-2d}
\end{equation}
At $\lambda_{xy}=0$, its represented Kronecker-sum discrepancy was zero and its
product-spectrum discrepancy was $3.55\times10^{-14}E_{\mathrm{ref}}$. At $\lambda_{xy}=0.30$, the
mixed-direction hopping norm was $1.94\times10^{-3}E_{\mathrm{ref}}$, compared with a
$3.00\times10^{-15}E_{\mathrm{ref}}$ numerical floor at zero coupling. On the strongest retained
coupling, the withheld RMS interpolation error decreased from
$5.75\times10^{-2}E_{\mathrm{ref}}$ for the onsite shell to $1.70\times10^{-6}E_{\mathrm{ref}}$ for the complete
finite transform. The latter reconstructed its own $15\times15$ mesh within
$1.11\times10^{-15}E_{\mathrm{ref}}$ while retaining nonzero error on an independent
$31\times31$ mesh.

These values are extracted from the earlier controlled scalar calculation. They
establish a historical finite-mesh separability, mixed-hopping, and off-node
interpolation control; they are not results of the new square-lattice rank-two protocol,
do not establish a 3-D outcome, and do not qualify a non-orthogonal lattice.

\section{Non-orthogonal but physically isotropic extension}

Let the direct-lattice basis be the columns of $A$ and define the direct-lattice metric
and reciprocal basis by
\begin{equation}
 G=A^{\mathsf T}A,
 \qquad B=2\pi A^{-\mathsf T}.
\end{equation}
For $\mathbf R\in\Z^d$, physical squared distance is
\begin{equation}
 r_G^2(\mathbf R)=\mathbf R^{\mathsf T}G\mathbf R.
 \label{eq:metric-radius}
\end{equation}
Thus the coordinate radius $R_x^2+R_y^2$ is not used for a non-orthogonal cell unless it
coincides with the declared metric. The integer lattice point group is
\begin{equation}
 \Gamma_G=\{M\in GL(d,\Z):M^{\mathsf T}GM=G\}.
 \label{eq:metric-group}
\end{equation}
An isotropic hopping inventory retains complete $\Gamma_G$ orbits and every attainable
metric shell through its cap. Equal coefficients are required only when imposed by the
model's point-group representation; support symmetry does not itself imply coefficient
equality.

A selected non-orthogonal parent must declare an inversion-complete and point-group-
complete bond orbit $\mathcal B$. A dimension-normalized candidate form is
\begin{equation}
 K_{\rho,G}(\mathbf q)=T_0+
 \frac{\rho}{|\mathcal B_+|}
 \sum_{\mathbf b\in\mathcal B_+}
 \left(e^{2\pi i\mathbf q\cdot\mathbf b}T_{\mathbf b}
 +e^{-2\pi i\mathbf q\cdot\mathbf b}T_{\mathbf b}^\dagger\right),
 \label{eq:general-parent}
\end{equation}
subject to the selected internal point-group representation. The smooth rotation must
likewise be constructed from a complete bond orbit rather than privileged coordinate
axes.

\subsection{Bond-orbit analytic comparator}

Let the chosen positive bond inventory be
$\mathcal B_+=\{\mathbf b_1,\ldots,\mathbf b_m\}\subset\Z^d\setminus\{\mathbf0\}$.
For real orbit weights $\eta_j$, set $W=\sum_j|\eta_j|>0$ and define the normalized
bond-orbit phase
\begin{equation}
 \theta_{a,\mathcal B}(\mathbf q)
 =\frac{a}{W}\sum_{j=1}^m\eta_j
   \sin(2\pi\mathbf q\cdot\mathbf b_j),
 \qquad |\theta_{a,\mathcal B}(\mathbf q)|\le a.
 \label{eq:bond-orbit-phase}
\end{equation}
The weights must be fixed by the selected symmetry contract; the theorem below does not
by itself declare them isotropic.

\begin{theorem}[Bond-orbit flat-parent comparator]
\label{thm:bond-orbit}
For $0\le a\le\pi/4$, define
\begin{equation}
 c_{\mathcal B,\boldsymbol\eta}(a)
 =\int_{\T^d}e^{2i\theta_{a,\mathcal B}(\mathbf q)}\,d\mathbf q.
 \label{eq:bond-orbit-correlation}
\end{equation}
Then $c_{\mathcal B,\boldsymbol\eta}(a)$ is real and strictly positive. If
\begin{align}
 D_{a,\mathcal B}(\mathbf q)
 &=R(\theta_{a,\mathcal B}(\mathbf q))^\dagger
   D R(\theta_{a,\mathcal B}(\mathbf q)),\\
 D_{\mathbf R}^{\mathcal B}
 &=\int_{\T^d}e^{-2\pi i\mathbf q\cdot\mathbf R}
   D_{a,\mathcal B}(\mathbf q)\,d\mathbf q,
\end{align}
then
\begin{align}
 D_{\mathbf0}^{\mathcal B}
 &=c_{\mathcal B,\boldsymbol\eta}(a)D,\\
 \inf_{s,\lambda\ge0,\phi}L_{\mathcal B,\phi}^2(s,\lambda)
 &=\kappa[1-c_{\mathcal B,\boldsymbol\eta}(a)^2]
 =\frac{1}{2E_{\mathrm{ref}}^2}\sum_{\mathbf R\ne\mathbf0}
   \norm{D_{\mathbf R}^{\mathcal B}}_F^2,
 \label{eq:bond-orbit-tail}\\
 \inf_\phi L_{\mathcal B,\phi}^2(0,1)
 &=2\kappa[1-c_{\mathcal B,\boldsymbol\eta}(a)],
\end{align}
where $L_{\mathcal B,\phi}$ is Eq.~\eqref{eq:flat-operator-loss} with
$B_a^{(d)}$ replaced by $D_{a,\mathcal B}$.

Moreover, the comparator has the exact zero-mode expansion
\begin{equation}
 c_{\mathcal B,\boldsymbol\eta}(a)
 =\sum_{\substack{\mathbf n\in\Z^m\\
                   \sum_{j=1}^m n_j\mathbf b_j=\mathbf0}}
   \prod_{j=1}^m
   J_{n_j}\!\left(\frac{2a\eta_j}{W}\right).
 \label{eq:bond-orbit-bessel-sum}
\end{equation}
If the bond vectors have no nontrivial integer relation, this reduces to
\begin{equation}
 c_{\mathcal B,\boldsymbol\eta}(a)
 =\prod_{j=1}^mJ_0\!\left(\frac{2a\eta_j}{W}\right).
 \label{eq:independent-bond-product}
\end{equation}
\end{theorem}

\noindent\emph{Proof.} See Appendix~\ref{app:proof-bond-orbit}.


\begin{theorem}[Constrained-series complete-shell tail certificate]
\label{thm:constrained-tail}
Let
\begin{equation}
 g_{\mathbf R}=
 \sum_{\substack{\mathbf n\in\mathbb Z^m\\
                   \sum_jn_j\mathbf b_j=\mathbf R}}
 \prod_{j=1}^mJ_{n_j}\!\left(\frac{2a\eta_j}{W}\right)
 \label{eq:constrained-coefficient}
\end{equation}
be the scalar Fourier coefficient of $e^{2i\theta_{a,\mathcal B}}$. In the constant
basis where $D=\delta\sigma_z$,
\begin{equation}
 D_{\mathbf R}^{\mathcal B}
 =\frac{\delta}{2}\left[
 (g_{\mathbf R}+g_{-\mathbf R})\sigma_z
 +i(g_{\mathbf R}-g_{-\mathbf R})\sigma_x
 \right].
 \label{eq:constrained-matrix-block}
\end{equation}
Let $\mathcal S_S$ be every lattice displacement in the complete physical-metric shells
through $S$. If finite constrained sums $\widetilde g_{\mathbf R}$ obey certified bounds
$|g_{\mathbf R}-\widetilde g_{\mathbf R}|\le\epsilon_{\mathbf R}$ for every
$\mathbf R\in\mathcal S_S$, then
\begin{equation}
 \frac{1}{2E_{\mathrm{ref}}^2}
 \sum_{\mathbf R\notin\mathcal S_S}
 \norm{D_{\mathbf R}^{\mathcal B}}_F^2
 \le\kappa\left\{
 1-\sum_{\mathbf R\in\mathcal S_S}
 \left[\max\{0,|\widetilde g_{\mathbf R}|-\epsilon_{\mathbf R}\}\right]^2
 \right\}.
 \label{eq:constrained-tail-certificate}
\end{equation}

For the selected equal-weight triangular orbit
$\{(1,0),(0,1),(1,-1)\}$, with $z=2a/3$,
\begin{equation}
 g_{(r_1,r_2)}=
 \sum_{\ell\in\mathbb Z}
 J_{r_1-\ell}(z)J_{r_2+\ell}(z)J_\ell(z).
 \label{eq:triangular-coefficient}
\end{equation}
For the selected equal-weight FCC orbit
\begin{equation}
 \{(1,0,0),(0,1,0),(0,0,1),(-1,1,0),(-1,0,1),(0,-1,1)\},
\end{equation}
with $z=2a/6$,
\begin{align}
 g_{(r_1,r_2,r_3)}
 =\sum_{u,v,w\in\mathbb Z}
 &J_{r_1+u+v}(z)J_{r_2-u+w}(z)J_{r_3-v-w}(z)\notag\\
 &\times J_u(z)J_v(z)J_w(z).
 \label{eq:fcc-coefficient}
\end{align}
Set $x=|z|/2$ and
\begin{align}
 A_0(z)&=e^{x^2}(2e^x-1),\\
 A_K(z)&=
 \frac{2e^{x^2}x^{K+1}}
 {(K+1)![1-x/(K+2)]}.
 \label{eq:absolute-bessel-tail}
\end{align}
Truncating Eq.~\eqref{eq:triangular-coefficient} to $|\ell|\le K$ gives
$\epsilon_{\mathbf R}=A_K(z)$. Truncating Eq.~\eqref{eq:fcc-coefficient} to
$|u|,|v|,|w|\le K$ gives
$\epsilon_{\mathbf R}=3A_K(z)A_0(z)^2$.
\end{theorem}

\noindent\emph{Proof.} See Appendix~\ref{app:proof-constrained-tail}.


For the selected maximum angle $a=\pi/4$, bounded binary64 evaluations with a declared
$10^{-13}$ per-coefficient numerical margin give normalized scalar-tail upper bounds
$2.8760\times10^{-9}$ for the triangular shell cap $S=21$ at $K=10$, and
$5.5358\times10^{-6}$ for the FCC shell cap $S=7$ at $K=9$. These values evaluate the
proved certificate in binary64 arithmetic with the stated margin; no interval-arithmetic
or mesh-limit extrapolation is applied.

The metric and point group select the physically admissible complete bond orbit; the
integer relations among the selected reduced-coordinate bond vectors determine whether
the comparator factorizes. Thus the square and cubic products arise from independent
primitive-axis harmonics, whereas a bond orbit with integer dependencies generally has
the constrained Bessel sum in Eq.~\eqref{eq:bond-orbit-bessel-sum} instead.

The non-orthogonal calculations use a graphene-inspired triangular/honeycomb 2-D
parent and a silicon-inspired FCC/diamond 3-D parent, together with matched Bravais
controls at equal reciprocal-torus RMS gauge amplitude. ``Inspired'' denotes only the
declared geometry, two-site basis, and bipartite connectivity; neither parent is a
material parametrization. The triangular and FCC control orbits yield, respectively,
one-index and three-index constrained Bessel sums under
Theorems~\ref{thm:bond-orbit} and \ref{thm:constrained-tail}. These dependent bond
orbits require their own metric shells and certificates; the square and cubic Bessel
products cannot be transferred unchanged.

\section{Restricted alignment and Hamiltonian comparison}

For $\mathcal U=\bigoplus_{\mathbf q}U(\mathbf q)$, define
\begin{equation}
 \widehat A_B=\mathcal U^\dagger\widehat H_B\mathcal U.
\end{equation}
The corresponding Wannier-space unitary
\begin{equation}
 \mathcal G_W=(\mathcal F_{d,N}^\dagger\otimes I_2)
 \mathcal U(\mathcal F_{d,N}\otimes I_2)
\end{equation}
is generally nonlocal and satisfies
\begin{equation}
 \widehat A_W=\mathcal G_W^\dagger\widehat H_W\mathcal G_W.
 \label{eq:covariant}
\end{equation}
The covariantly aligned subtraction vanishes by construction. The restricted scientific
comparison replaces $\mathcal G_W$ by one cell-local global rotation
$I_{N^d}\otimes R(\phi)$. A parent-space comparison requires the retained projector and
an explicit lift. Every subtraction must state its source and target spaces, basis,
gauge, energy reference, units, geometry, and normalization.

\section{Synthetic-strain methods}
\label{sec:strain-protocol}

\subsection{Scientific question and three controls}

The strain protocol asks how a homogeneous change of lattice metric affects the
mechanisms and sampled requirement coordinates associated with spectral fidelity,
restricted operator agreement, and finite-range hopping locality. It keeps retained
rank, energy reference, and reduced-coordinate gauge explicit. It is a synthetic gauge
and model-response stress test, not an elastic, structural-stability, or material-strain
calculation. The deformation does not make the exact spectral and hopping
representations inconsistent. The fixed-geometry Pareto surfaces and strain-response
scans remain separate until explicit strain-resolved landscapes are calculated.

The protocol separates three controls:
\begin{enumerate}
 \item The \emph{geometry-only control} uses a constant rank-two parent and an imposed
 smooth reduced-coordinate gauge. The Hamiltonian field, Fourier coefficients, exact
 band equivalence, and covariant residual are strain invariant; only the physical metric
 used to group complete hopping shells changes.
 \item The \emph{frozen-parent null-response control} uses the four-state
 bipartite parent, fixed onsite and hopping matrices, the seed-projected retained-frame
 construction, and the explicit common-cell map. A nonzero aligned strain response is a
 failed null test rather than a material signal.
 \item The \emph{synthetic radial-response control} starts from the same parent and
 onsite matrices but changes only selected inter-sublattice bond amplitudes according to
 \begin{equation}
  t_j(s)=t_j(0)\exp\!\left[-\beta
  \left(\frac{d_j(s)}{d_j(0)}-1\right)\right],
  \qquad \beta=1.
  \label{eq:strain-radial-response-law}
 \end{equation}
 No onsite response, strain-dependent refit, normalization, structural relaxation, or
 averaging with either other control is permitted.
\end{enumerate}
The three controls answer different questions and are not successive approximations to
one Hamiltonian.

\subsection{Shell--scale Pareto decision landscape}

At strain coordinate $s$, let $H_s(\mathbf q)$ be the retained reference field,
$\bar h_s(\mathbf q)=\operatorname{tr}H_s(\mathbf q)/2$, and
$\widetilde H_s(\mathbf q)=H_s(\mathbf q)-\bar h_s(\mathbf q) I_2$. Let
$B_{s,S}(\mathbf q)$ be the gauge-transformed
comparison reconstructed after retaining every complete physical-metric shell through
$S$. The restricted full-field alignment supplies one global rotation $R_{*,s}$ and a
nonnegative optimized scale $\lambda_{*,s}$. We sample the scientifically interpretable
path
\begin{equation}
 \lambda_s(u)=(1-u)\lambda_{*,s}+u,
 \qquad 0\le u\le1,
 \label{eq:scale-path}
\end{equation}
so that $u=0$ is the optimized operator-scale endpoint and $u=1$ preserves the original
spectral splitting. The corresponding globally aligned target is
\begin{equation}
 A_{s,u}(\mathbf q)=\bar h_s(\mathbf q)I_2+
 \lambda_s(u)R_{*,s}^{\mathsf T}\widetilde H_s(\mathbf q)R_{*,s}.
 \label{eq:scale-path-target}
\end{equation}
For each calculated pair $(S,u)$, define
\begin{align}
 \varepsilon_H(s;S,u)
 &=\max_{\mathbf q}\frac{\norm{B_{s,S}(\mathbf q)-A_{s,u}(\mathbf q)}_F}
 {\sqrt2E_{\mathrm{ref}}},\\
 \varepsilon_E(s;S,u)
 &=\max_{\mathbf q,n}\frac{
 \left|\varepsilon_n[B_{s,S}(\mathbf q)]-
 \varepsilon_n[A_{s,u}(\mathbf q)]\right|}{E_{\mathrm{ref}}},\\
 \mathcal J_s(S,u)
 &=\left(N_s(S),\varepsilon_H(s;S,u),\varepsilon_E(s;S,u)\right).
 \label{eq:pareto-decision-objectives}
\end{align}
Lower values of all three components are preferred. A calculated point is Pareto
efficient when no other sampled pair weakly improves every component and strictly
improves at least one. The strain coordinate therefore indexes a family of sampled
Pareto decision landscapes; it is not itself minimized. The physical shell radius
$R_s(S)$ remains a separately reported locality coordinate because it need not order
shells in the same way as $N_s(S)$.

This construction requires no acceptance tolerance. A tolerance is needed only later to
select an operating point for a specified application or material. Tessellated cells may
connect adjacent calculated $(S,u)$ samples for visualization, but they assert neither
continuum coverage nor interpolation validity.

\subsection{Area- and volume-preserving paths}

We use the following area- and volume-preserving crystallographic paths and coordinate
bridges.

Let $s\in[0,1]$ be the dimensionless strain fraction. The 2-D unit-area rhombic path is
\begin{align}
 \theta(s)&=\frac\pi3+s\frac\pi6,\\
 A_2(s)&=\frac{1}{\sqrt{\sin\theta(s)}}
 \begin{pmatrix}
 1&\cos\theta(s)\\
 0&\sin\theta(s)
 \end{pmatrix},
 \qquad \det A_2(s)=1.
 \label{eq:strain-rhombic-path}
\end{align}
It connects triangular metric geometry at $s=0$ to square metric geometry at $s=1$.

For the 3-D Bain path, write $t=s\in[0,1]$ and define the tetragonal aspect ratio
and volume-compensating transverse and longitudinal scales by
\begin{equation}
 \gamma(t)=\sqrt2+t(1-\sqrt2),
 \qquad
 \ell_\perp(t)=\left(\frac{2}{\gamma(t)}\right)^{1/3},
 \qquad
 \ell_\parallel(t)=\gamma(t)\ell_\perp(t).
 \label{eq:strain-bain-scales}
\end{equation}
The three primitive vectors evolve continuously as
\begin{align}
 \mathbf a_1(t)&=\frac12
 \begin{pmatrix}\ell_\perp(t)\\\ell_\perp(t)\\\ell_\parallel(t)\end{pmatrix},
 &0\le t\le1,\\
 \mathbf a_2(t)&=\frac12
 \begin{pmatrix}-\ell_\perp(t)\\\ell_\perp(t)\\\ell_\parallel(t)\end{pmatrix},
 &0\le t\le1,\\
 \mathbf a_3(t)&=\frac12
 \begin{pmatrix}\ell_\perp(t)\\-\ell_\perp(t)\\\ell_\parallel(t)\end{pmatrix},
 &0\le t\le1.
 \label{eq:strain-bain-vectors}
\end{align}
Equivalently,
$A_3(t)=[\mathbf a_1(t)\ \mathbf a_2(t)\ \mathbf a_3(t)]$. Its primitive-cell volume
is exactly
\begin{equation}
 V(t)=\left|\mathbf a_1(t)\cdot
 [\mathbf a_2(t)\times\mathbf a_3(t)]\right|
 =\det A_3(t)
 =\frac{\ell_\perp(t)^2\ell_\parallel(t)}{2}=1
 \qquad (0\le t\le1).
 \label{eq:strain-bain-volume-invariant}
\end{equation}
Thus the complete path, including its endpoints, is volume preserving. The ratio
$\ell_\parallel/\ell_\perp=\gamma$ runs from $\sqrt2$ at $t=0$ to $1$ at $t=1$,
giving the FCC and BCC endpoint representations of the Bain path
\cite{Bain1924Martensite}.

\subsection{FCC--Bain common-cell identification}

Let $\mathbf n_F$ denote the existing FCC primitive coordinate and $\mathbf n_B$ the
Bain-coordinate label. The fixed unimodular map is
\begin{equation}
 \mathbf n_F=U\mathbf n_B,
 \qquad
 U=\begin{pmatrix}
 1&1&0\\
 0&0&1\\
 0&-1&0
 \end{pmatrix},
 \qquad \det U=1.
 \label{eq:strain-common-cell-map}
\end{equation}
It induces
\begin{equation}
 \mathbf q_B=U^{\mathsf T}\mathbf q_F,
 \qquad
 \mathbf n_B\longmapsto U\mathbf n_B\pmod N
 \label{eq:strain-reciprocal-cell-map}
\end{equation}
on the reciprocal torus and every finite periodic cell mesh. The diamond basis offset
$(1/4,1/4,1/4)_F$ becomes $(1/2,-1/4,1/4)_B$; orbital and sublattice labels are
unchanged. At the endpoint, $G_B=U^{\mathsf T}G_FU$, and the six FCC bond pairs map to
the six Bain-coordinate pairs. These identities establish a coordinate and state-space
map only; they do not identify a material Hamiltonian across strain.

\subsection{Sampling and numerical diagnostics}

The numerical protocol samples
\begin{equation}
 s\in\{0,1/8,2/8,\ldots,1\},
 \qquad a=\pi/4,
 \qquad \tau_{\mathrm{tail}}\in\{10^{-3},10^{-5}\}.
 \label{eq:strain-frozen-samples}
\end{equation}
The geometry-only transform uses $64^2$ points in 2-D and $24^3$ points in 3-D. The
frozen-parent and radial-response controls use construction/evaluation sizes
$16^2/33^2$ in 2-D and $8^3/17^3$ in 3-D, with evaluation shift $1/3$. Both response
controls fix $\rho=1$, the onsite matrices, $E_{\mathrm{ref}}$, and the imposed
reduced-coordinate gauge at every strain point. These finite meshes define the reported
synthetic scans; they are not extrapolated to a continuum-mesh limit.

At every point the protocol reports the unit-determinant metric, bond-length extrema and
orbit splitting, covariant and restricted-alignment meanings, ordinary and centered
Nyquist Fourier power, and the first complete physical-radius support meeting each tail
tolerance. Equal-radius ties are never split. The two response controls additionally use
the matrix and sorted-band interpolation, deliberate-truncation, and total channels
defined in Appendix~\ref{app:diagnostics}, both for the Hamiltonian and for the aligned
strain difference $\Delta H_s$. The three error channels are not added.

The bounded scan uses the base grid
\begin{equation}
 \widehat S_0=\{s_i=i/200:i=0,\ldots,200\}.
 \label{eq:strain-scan-grid}
\end{equation}
A hat denotes a maximum over these samples, not a continuum supremum. If $\kappa_i$ is
the retained splitting scale and $\mu_i$ the maximum finite-mesh correlation after one
global real rotation, the pointwise optimized operator requirement and the requirement
for the unchanged reference candidate $p_0=(0,1)$ are
\begin{align}
 \tau_{o,\min}(s_i)
 &=\sqrt{\kappa_i[1-\max(0,\mu_i)^2]},\\
 \tau_{o,p_0}(s_i)
 &=\sqrt{2\kappa_i(1-\mu_i)}.
 \label{eq:strain-scan-witness-requirements}
\end{align}
Their sampled envelopes are $\widehat\tau_{o,\min}=\max_i\tau_{o,\min}(s_i)$ and
$\widehat\tau_{o,p_0}=\max_i\tau_{o,p_0}(s_i)$. Pointwise optimized candidates are not
one common path witness, and $p_0$ is only a prescribed candidate path until frozen
limits establish feasibility. Likewise,
$\widehat R_\varepsilon=\max_iR_\varepsilon(s_i)$ and
$\widehat N_\varepsilon=\max_iN_\varepsilon(s_i)$ are reported separately because their
maxima need not occur at the same strain.

Every adjacent base-grid interval whose endpoint block counts or bond-degeneracy counts
differ is bisected five times, yielding detected transition brackets of width at most
$1/6400$. This bounded rule does not detect an even number of hidden changes inside an
equal-endpoint interval. It supplies sampled localization only, not a continuum
certificate. No feasibility crossing is assigned: this protocol characterizes the sampled decision
landscape rather than classifying it against an acceptance region.

The calculations are synthetic and do not determine forces, stresses, strain energies,
phase stability, realistic graphene or silicon response, or a material FCC--BCC
transition. Material-specific strained Hamiltonians and environment-responsive models
such as Refs.~\cite{Fang2018Strained2D,Zhang2022EquivariantHamiltonians} remain outside
the present scope.

\section{Diagnostic summary}

The study keeps parent and retained-space reconstruction, covariant equivalence,
restricted global alignment, cyclic conversion, interpolation, deliberate truncation,
total matrix and band consequences, locality power, transport, and refinement as
distinct diagnostic layers. Appendix~\ref{app:diagnostics} defines their
normalizations, applicability conditions, and first-passing rules. In particular, the
three approximation channels in Eq.~\eqref{eq:three-error-channels} are not additive,
and the closure-phase reduction in
Eq.~\eqref{eq:multidimensional-closure-phase-reduction} is defined only when every
contributing transport link is full rank.

Appendix~\ref{app:verification} summarizes the independent numerical reconstruction.

\section{Results}

The calculations use three construction levels, $16^2$, $32^2$, and $64^2$ in two
dimensions and $8^3$, $16^3$, and $32^3$ in three dimensions, together with shifted
evaluation meshes. Reconstruction, Hermiticity, Parseval, transport, and closure
identities remain below $2.0\times10^{-14}$. One of 18 graphene-family first-passing
shell decisions changes by one shell between the first two levels; every such decision
is unchanged on the final refinement in all four families. Some deliberately strict
tail cuts remain outside the finite shell caps, so the central comparison uses the full
sampled decision landscapes rather than converting those cuts into acceptance tests.

\paragraph{Central result: controlled incompatibility landscape.}
The exact untruncated route preserves the retained spectrum under the declared
momentum-dependent gauge. Once the comparison is restricted to one cell-local global
rotation, one scale choice, and a finite complete-shell inventory, spectral fidelity,
Hamiltonian fidelity, and locality cannot in general be optimized simultaneously.
Shell cutoff controls retained complexity, while the path from $\lambda_{*,s}$ to unit
scale resolves the compromise between restricted operator agreement and preservation of
the original splitting. Synthetic deformation is analyzed through separate response
channels, while the reported Pareto surfaces remain fixed-geometry results. Exact gauge
equivalence holds throughout both calculations.

The geometry-only control isolates metric reordering of the same Fourier inventory; the
frozen-parent control tests that coordinate deformation alone produces no Hamiltonian
response; and the radial-response control deliberately changes bond amplitudes. Taken
together, the controls separate mechanisms rather than plotting an undifferentiated
strain response. The development output here is the set of sampled nondominated
compromises, not a self-selected acceptance boundary. A componentwise cover assembled
from different samples is not an attainable decision and is therefore not a Pareto
representative.

The bounded fixed-geometry calculation samples 17 scales between $\lambda_*$ and unity
at every available complete shell for each of the eight cases. The optimized scale rises
from $0.729890$ for the square directed parent to $0.937745$ for the diamond-inspired
case. Every case has a nonempty, non-collapsed nondominated set, but its sampled shape is
case dependent; this calculation does not isolate metric-shell geometry from parent
connectivity. The diamond-inspired case is the most selective in this sampling,
with 34 of 136 calculated vertices nondominated; the triangular and FCC controls retain
156 of 187 and 116 of 136, respectively. At the largest sampled diamond shell, moving
from $\lambda_*$ to unit scale increases the normalized Hamiltonian mismatch from
$0.211780$ to $0.217933$ (about $2.9\%$) while reducing the sorted-band mismatch from
$0.0193798$ to $0.00341626$ (about $82\%$). This is one endpoint comparison, not a
universal preference for unit scale. The vertex counts describe the declared discrete
grid and are not invariant geometric measures. Figure~\ref{fig:principal-decision-landscapes}
shows the four principal cases; Appendix~\ref{app:case-decision-landscapes} gives all
eight cases separately.

\begin{figure}[p]
 \centering
 \includegraphics[width=0.98\textwidth]{figures/principal-shell-scale-decision-landscapes.pdf}
 \caption{Sampled shell--scale decision landscapes for the four principal cases. Each
 surface uses every complete shell and 17 values along the path from the optimized
 positive comparison scale to unit scale. Orange markers are nondominated calculated
 vertices for the three minimized objectives: retained blocks, maximum Hamiltonian
 error, and maximum sorted-band error. Tessellation connects adjacent calculated
 shell--scale samples and represents the discrete sampling topology.}
 \label{fig:principal-decision-landscapes}
\end{figure}

Figures~\ref{fig:family-square}--\ref{fig:family-silicon-inspired} compare the four
model families at fixed geometry; the strain-responsive controls are reported
separately.

\begin{figure}[p]
 \centering
 \includegraphics[width=0.98\textwidth]{figures/square-2d-scientific-story.pdf}
 \caption{Square-lattice results. The panels separate withheld interpolation,
 external retained-subspace isolation, hopping-tail power, withheld spectral
 consequence, and the relation between tail power and spectral consequence.}
 \label{fig:family-square}
\end{figure}

\begin{figure}[p]
 \centering
 \includegraphics[width=0.98\textwidth]{figures/triangular-honeycomb-2d-scientific-story.pdf}
 \caption{Graphene-inspired honeycomb results and matched Bravais controls. The
 non-orthogonal physical metric changes the complete-shell inventory relative to the
 square family.}
 \label{fig:family-graphene-inspired}
\end{figure}

\begin{figure}[p]
 \centering
 \includegraphics[width=0.98\textwidth]{figures/simple-cubic-3d-scientific-story.pdf}
 \caption{Simple-cubic results at two displayed mesh levels. The separate panels retain
 interpolation, external isolation, hopping-tail power, and spectral consequence as
 distinct quantities.}
 \label{fig:family-simple-cubic}
\end{figure}

\begin{figure}[p]
 \centering
 \includegraphics[width=0.98\textwidth]{figures/fcc-diamond-3d-scientific-story.pdf}
 \caption{Silicon-inspired diamond results and matched FCC Bravais controls. The two
 cases share the FCC metric but differ in parent connectivity.}
 \label{fig:family-silicon-inspired}
\end{figure}

Figures~\ref{fig:three-channel-square}--\ref{fig:three-channel-silicon-inspired}
show the corresponding refined-mesh three-channel decompositions. Near overlap of the
deliberate-truncation and total curves means interpolation is much smaller at these
sampled points; it does not make the channels definitionally identical.

\begin{figure}[p]
 \centering
 \includegraphics[width=0.98\textwidth]{figures/square-2d-three-channel-error-decomposition.pdf}
 \caption{Square-lattice endpoint comparison: maximum Hamiltonian and sorted-band
 interpolation, deliberate-truncation, and total error channels.}
 \label{fig:three-channel-square}
\end{figure}

\begin{figure}[p]
 \centering
 \includegraphics[width=0.98\textwidth]{figures/triangular-honeycomb-2d-three-channel-error-decomposition.pdf}
 \caption{Graphene-inspired family and matched controls: refined-mesh endpoint
 three-channel decomposition. Every row retains its own complete physical-metric shell
 inventory.}
 \label{fig:three-channel-graphene-inspired}
\end{figure}

\begin{figure}[p]
 \centering
 \includegraphics[width=0.98\textwidth]{figures/simple-cubic-3d-three-channel-error-decomposition.pdf}
 \caption{Simple-cubic endpoint comparison: maximum Hamiltonian and sorted-band
 interpolation, deliberate-truncation, and total error channels.}
 \label{fig:three-channel-simple-cubic}
\end{figure}

\begin{figure}[p]
 \centering
 \includegraphics[width=0.98\textwidth]{figures/fcc-diamond-3d-three-channel-error-decomposition.pdf}
 \caption{Silicon-inspired family and matched controls: refined-mesh endpoint
 three-channel decomposition. These are errors for the declared synthetic model, not
 silicon material errors.}
 \label{fig:three-channel-silicon-inspired}
\end{figure}

\paragraph{Bounded geometry-only control.}
Separate frozen-parent geometry-stress diagnostics use nine unit-area rhombic and
unit-volume Bain metrics. The reduced-coordinate Hamiltonian, gauge, Fourier-block
weights, and restricted-alignment residual are invariant by construction. Nevertheless,
the changing physical metric reorders complete hopping shells and changes the sampled
first-passing physical radius and retained block count. At tail-power tolerance
$10^{-5}$, the endpoint radius changes from $3.7224$ to $4.2426$ in 2-D and from
$2.9698$ to $3.2734$ in 3-D; the corresponding block counts change from $43$ to $61$
and from $135$ to $169$. At the looser 3-D tolerance $10^{-3}$, the endpoint radius
instead changes from $2.2449$ to $2.1822$ while the block count changes from $55$ to
$59$. Physical range and parameter count are therefore distinct locality costs even
when the abstract reduced-coordinate operator is unchanged. A unimodular map identifies
the existing FCC primitive coordinates with the Bain endpoint coordinates.
Figures~\ref{fig:geometry-control-strain-2d} and
\ref{fig:geometry-control-strain-3d} display the complete bounded geometry-only
diagnostics.

\begin{figure}[p]
 \centering
 \includegraphics[width=0.98\textwidth]{figures/rhombic-2d-geometry-only-control.pdf}
 \caption{Two-dimensional frozen-gauge geometry-only result. The parent, gauge,
 Fourier weights, and unit-cell area remain fixed while the rhombic metric moves from
 triangular to square. Panels separate bond-orbit splitting, invariant restricted
 residual, first-passing physical radius, and retained block count. The plotted
 tail-power thresholds are locality diagnostics, not universal spectral-error bounds.}
 \label{fig:geometry-control-strain-2d}
\end{figure}

\begin{figure}[p]
 \centering
 \includegraphics[width=0.98\textwidth]{figures/bain-3d-geometry-only-control.pdf}
 \caption{Three-dimensional frozen-gauge geometry-only result along the unit-volume
 FCC--BCC Bain metric path. The looser threshold shows that first-passing physical
 radius and retained block count follow different trends along the same metric path.}
 \label{fig:geometry-control-strain-3d}
\end{figure}

\paragraph{Bounded frozen-parent and radial-response controls.}
The response calculation uses the four-state bipartite parent. The frozen-parent
null-response control holds its onsite and hopping matrices
fixed; its pointwise gauge-aligned strain difference remains below
$1.4\times10^{-15}E_{\mathrm{ref}}$ on both paths. The synthetic radial-response
control changes only the parent-bond amplitudes through the prescribed
dimensionless radial law
$t_j(s)/t_j(0)=\exp[-(d_j(s)/d_j(0)-1)]$. At the final 2-D metric, the responses span
$0.81758$--$1.27146$ and the maximum aligned retained-space perturbation is
$0.063575E_{\mathrm{ref}}$. At the BCC endpoint, the four selected bonds share response
$0.97564$ and the corresponding maximum is $0.007289E_{\mathrm{ref}}$. The stored
trace/traceless decomposition is dominated by the traceless channel on both paths. At
the $10^{-5}$ tail support, the 2-D endpoint interpolation, deliberate-
truncation, and total maximum matrix errors of $\Delta H_s$ are approximately
$6.0\times10^{-8}$, $7.801\times10^{-4}$, and $7.802\times10^{-4}$; the 3-D values
are $9.36\times10^{-5}$, $1.61\times10^{-4}$, and $1.89\times10^{-4}$. These
channels are displayed in Figs.~\ref{fig:response-control-strain-2d} and
\ref{fig:response-control-strain-3d} and are not added.

\begin{figure}[p]
 \centering
 \includegraphics[width=0.98\textwidth]{figures/rhombic-2d-strain-response-controls.pdf}
 \caption{Two-dimensional response controls. The frozen-parent calculation gives the
 null aligned difference expected from the unchanged parent. The radial-response
 control uses the fixed synthetic law and reports its aligned perturbation,
 scalar/traceless split, locality cost, and three approximation channels.}
 \label{fig:response-control-strain-2d}
\end{figure}

\begin{figure}[p]
 \centering
 \includegraphics[width=0.98\textwidth]{figures/bain-3d-strain-response-controls.pdf}
 \caption{Three-dimensional Bain-path response controls. The prescribed radial law
 generates a finite response while the frozen parent remains a numerical null.}
 \label{fig:response-control-strain-3d}
\end{figure}

\paragraph{Pathwise qualification and approximation envelopes.}
The 201-point response sweeps retain the external-gap, seed-overlap, frame-conditioning,
approximation, and mechanism extrema together with the first strain sample where each
occurs. For the radial-response control, the minimum sampled external isolation gaps are
$2.00708E_{\mathrm{ref}}$ in 2-D and $2.01305E_{\mathrm{ref}}$ in 3-D; the minimum
seed overlaps are $0.998684$ and $0.998687$. The minimum frame-overlap singular values
remain above $0.9999998$ and $0.999999996$, respectively. Thus the sampled response does
not coincide with loss of the retained subspace or a rank-deficient frame overlap.

The maximum radial-response pointwise frame-alignment residuals are
$8.27192\times10^{-4}$ in 2-D and $9.64819\times10^{-5}$ in 3-D, both at $s=1$.
At the $10^{-5}$ support, the sampled maxima of the total aligned-difference Hamiltonian
error are $7.80164\times10^{-4}E_{\mathrm{ref}}$ in 2-D and
$1.89367\times10^{-4}E_{\mathrm{ref}}$ in 3-D. The corresponding frozen-parent
envelopes remain at binary64 roundoff. In both radial controls the response is
traceless-dominated: its maximum scalar/traceless components are
$1.77497\times10^{-6}/0.0635748$ in 2-D and
$8.93927\times10^{-7}/0.00728904$ in 3-D. These are maxima over calculated samples,
not continuum bounds. Figures~\ref{fig:qualification-response-scan-2d} and
\ref{fig:qualification-response-scan-3d} retain the corresponding pointwise envelopes.

\begin{figure}[p]
 \centering
 \includegraphics[width=0.98\textwidth]{figures/rhombic-2d-strain-qualification-response-scan.pdf}
 \caption{Two-dimensional 201-point qualification and response-mechanism scan. External
 gap, seed overlap, frame conditioning, radial-response magnitude and mechanism split,
 and tight-support approximation errors remain distinct channels.}
 \label{fig:qualification-response-scan-2d}
\end{figure}

\begin{figure}[p]
 \centering
 \includegraphics[width=0.98\textwidth]{figures/bain-3d-strain-qualification-response-scan.pdf}
 \caption{Three-dimensional 201-point Bain qualification and response-mechanism scan. The
 synthetic diamond radial response remains separated from retained-subspace and frame
 qualification throughout the sampled path.}
 \label{fig:qualification-response-scan-3d}
\end{figure}

\paragraph{Sampled witness and transition scan.}
The 201-point base scan preserves the normalized cell measure to
$4.5\times10^{-16}$ in 2-D and $1.3\times10^{-15}$ in 3-D; direct sorted-band defects
remain below $1.8\times10^{-15}$. In 2-D the synthetic radial response increases the
pointwise optimized requirement from $0.0810724$ at $s=0$ to $0.0875414$ at $s=1$ and
the fixed-$p_0$ requirement from $0.0835420$ to $0.0907685$. In 3-D both radial-response
requirements attain their sampled maxima at $s=0$ and decrease slightly toward the BCC
endpoint. The geometry-only and frozen-parent requirement curves are strain invariant,
as their control definitions demand. These are sampled requirement coordinates within the decision landscape, not pass/fail
results; no acceptance classification is intended for the synthetic controls.

The sampled 2-D locality envelopes are $(\widehat R,\widehat N)=(2.82843,25)$ at
$10^{-3}$ and $(4.24264,61)$ at $10^{-5}$, all attained at $s=1$. In 3-D, the looser
tolerance separates the two costs: $\widehat R=2.24492$ occurs at $s=0$, whereas
$\widehat N=59$ first occurs at $s=0.77$; the tighter envelope is
$(3.27337,169)$ at $s=1$. Five bisections localize seven detected 2-D shell-count
changes and six detected 3-D changes to intervals no wider than $1/6400$. On both paths
the endpoint bond-degeneracy change is bracketed by $[0,1/6400]$. These brackets do not
exclude hidden equal-endpoint transitions. Independent test-only finite sums reconstruct
the first-passing shells, and an independent harmonic calculation reconstructs both
operator-requirement channels. Figures~\ref{fig:strain-scan-2d} and
\ref{fig:strain-scan-3d} show the sampled envelopes and controls.

\begin{figure}[p]
 \centering
 \includegraphics[width=0.98\textwidth]{figures/rhombic-2d-strain-scan.pdf}
 \caption{Two-dimensional 201-point rhombic scan. Pointwise and fixed-candidate
 requirements, geometry-only locality costs, the frozen-parent null response, and the
 synthetic radial response remain separate across all 201 sampled coordinates.}
 \label{fig:strain-scan-2d}
\end{figure}

\begin{figure}[p]
 \centering
 \includegraphics[width=0.98\textwidth]{figures/bain-3d-strain-scan.pdf}
 \caption{Three-dimensional 201-point Bain scan. At the looser tail tolerance, the sampled
 maximum physical radius and maximum block count occur at different sampled strains.}
 \label{fig:strain-scan-3d}
\end{figure}

Within these replaceable finite sampled development calculations, metric deformation
reorganizes complete physical shells, an unchanged parent produces no resolved
Hamiltonian response, and the prescribed radial law produces a finite, predominantly
traceless response without loss of retained-subspace isolation. They are not qualified
retained evidence and do not define a material strain potential or establish a continuum
extremum.

\clearpage
\section{Discussion}

\subsection{The three-way decision landscape}

The sampled development surfaces show that spectral fidelity, Hamiltonian fidelity, and
hopping locality form a three-way trade-off once two practical restrictions are imposed:
alignment is limited to a single cell-local rotation, and hopping support is limited to
a finite set of complete metric shells. Exact gauge equivalence continues to hold in
the unrestricted setting; the incompatibility appears only after those restrictions are
added.

Consequently, a small sorted-band error is not by itself sufficient to establish the
operator fidelity or locality of a finite-range Hamiltonian representation. The sampled
Pareto fronts for all eight cases are nonempty and non-collapsed. In every case, some
calculated points improve locality only by degrading operator or spectral agreement, and
vice versa. The development output is therefore the set of sampled nondominated
compromises rather than a single preferred Hamiltonian.

\subsection{Case dependence and the role of the comparison scale}

The shape of the decision surface is strongly case dependent. The present comparison
does not isolate lattice metric and shell geometry from parent connectivity. The
diamond-inspired case is the most selective of the eight fronts examined: only 34 of 136 calculated
vertices are nondominated, whereas the triangular and FCC controls retain substantially
larger fractions of their sampled vertices. Similar retained-block counts can correspond
to different physical ranges once the metric changes, and equal shell indices have no
intrinsic cross-lattice meaning without the associated metric and orbit inventory.

The path from the optimized scale $\lambda_*$ to unit scale quantifies a concrete
compromise. At the largest sampled diamond shell, restoring unit scale increases the
normalized Hamiltonian mismatch by about $2.9\%$ while reducing the sorted-band mismatch
by approximately $82\%$. Neither endpoint is universally preferred: $\lambda_*$ favors
restricted operator agreement, whereas unit scale recovers most of the original
spectral splitting at modest additional operator cost. The appropriate operating point
is application dependent and requires an external criterion.

\subsection{Mechanism separation under synthetic strain}

The three synthetic controls isolate distinct contributions to the observed trade-offs.
Geometry-only deformation leaves the reduced-coordinate Hamiltonian, gauge, and Fourier
weights invariant, yet reorders complete physical shells and changes both first-passing
radius and retained block count. At the looser three-dimensional tail tolerance, these
two locality costs move in opposite directions. Physical range and parameter count are
therefore independent locality costs.

The frozen-parent control returns a numerical null, with aligned differences at
binary64 roundoff. Coordinate deformation alone produces no Hamiltonian response when
the parent matrices are fixed. The prescribed radial-response control produces a finite,
predominantly traceless perturbation while preserving external isolation of the retained
subspace and full-rank frame conditioning. A finite response can therefore be introduced
without losing the spectral isolation required by the representation chain.

\subsection{Implications and outlook}

The rank-two benchmark supplies a controlled setting in which finite-range models are
assessed against an explicit three-objective surface rather than a single error metric.
In Wannier or machine-learned tight-binding workflows, an application can impose its own
Hamiltonian and spectral tolerances; the nondominated set then identifies the associated
locality choices.

Several extensions follow directly. Raising the retained rank above two will test how
the residual--locality identities and Pareto-front structure change beyond the simple
$U(2)$ gauge freedom used here. Direction-dependent coupling laws and material-specific
parent Hamiltonians can enter the same diagnostic channels without changing the
underlying state-space identifications. Repeating the shell--scale calculation along a
denser strain path or a lattice-relaxation trajectory would resolve the corresponding
motion of the decision landscape.

The essential distinction remains: exact spectral equivalence alone is insufficient;
restricted alignment and finite hopping support introduce the nontrivial representation
cost.

\section{Conclusions}

We have separated exact spectral equivalence from the additional demands required of a
finite-range tight-binding representation. Spectral projection can generate infinite
support, and restricted global alignment exposes a residual that is exactly tied to
non-onsite Fourier weight in the orthogonal flat-parent family. Metric-complete shells
and bond-orbit identities extend the construction to non-orthogonal isotropic lattices.
Across eight calculated cases, locality, Hamiltonian fidelity, and sorted-band fidelity
form nontrivial sampled Pareto sets rather than sharing one optimum. Synthetic rhombic
and Bain controls further separate geometric reordering, a frozen-parent null response,
and a finite radial response. Taken as replaceable development diagnostics rather than
qualified retained evidence, these results define a controlled rank-two benchmark for
assessing reduced Hamiltonians while preserving the distinctions among gauge
equivalence, operator agreement, spectral agreement, and real-space support.

\clearpage
\appendix
\section{Proofs}
\label{app:proofs}

The theorem and corollary statements remain in the main text. Their proofs are collected
here so that the main narrative can emphasize the representation chain, geometric
controls, and numerical evidence.

\subsection{Spectral projection and finite hopping support}
\label{app:proof-projection-range}

\begin{proof}[Proof of Proposition~\ref{prop:projection-destroys-range}]
Set $z=e^{ik}$. If $d(k)$ were a finite Fourier polynomial, there would be a Laurent
polynomial $f(z)$ whose restriction to the unit circle equals $d(k)$. Squaring on the
unit circle and using equality on infinitely many points would give the Laurent-polynomial
identity
\begin{equation}
 f(z)^2=\Delta^2+2t^2+t^2(z+z^{-1}).
 \label{eq:laurent-square}
\end{equation}
If $m$ and $n$ are the smallest and largest exponents with nonzero coefficients in
$f$, then the smallest and largest exponents in $f^2$ are $2m$ and $2n$; their
coefficients cannot cancel because each is the square of an extreme coefficient. The
right-hand side of Eq.~\eqref{eq:laurent-square} has extreme exponents $-1$ and $1$
when $t\ne0$, which is impossible because $2m$ and $2n$ are even. Thus $d$ is not a
finite Fourier polynomial. Equation~\eqref{eq:projected-range} has diagonal entries
containing $-d(k)/2$ plus finite Fourier terms, so the embedded projected operator is not
finite range.
\end{proof}

\subsection{Flat-parent residual--locality identity}
\label{app:proof-flat-parent}

\begin{proof}[Proof of Theorem~\ref{thm:flat-parent}]
A constant real rotation diagonalizes $D$ without changing any Frobenius norm, Fourier
weight, or minimization in the statement, so take $D=\delta\sigma_z$ with
$\delta\ne0$. Conjugation through a real rotation gives
\begin{equation}
 B_a^{(d)}(\mathbf q)
 =\delta\left[
 \cos\!\left(2\theta_a^{(d)}(\mathbf q)\right)\sigma_z
 -\sin\!\left(2\theta_a^{(d)}(\mathbf q)\right)\sigma_x
 \right].
 \label{eq:flat-rotated-pauli}
\end{equation}
Using the normalized integral representation
\begin{equation}
 J_0(x)=\int_{-1/2}^{1/2}e^{ix\sin(2\pi q)}\,dq
 \label{eq:bessel-integral}
\end{equation}
and the product structure of $\T^d$,
\begin{align}
 \int_{\T^d}e^{2i\theta_a^{(d)}(\mathbf q)}\,d\mathbf q
 &=\prod_{\alpha=1}^d
   \int_{-1/2}^{1/2}
   e^{i(2a/d)\sin(2\pi q_\alpha)}\,dq_\alpha\\
 &=J_0\!\left(\frac{2a}{d}\right)^d=c_d(a).
\end{align}
Global inversion $\mathbf q\mapsto-\mathbf q$ makes the sine average in
Eq.~\eqref{eq:flat-rotated-pauli} vanish, proving
Eq.~\eqref{eq:onsite-average}. Moreover,
$0\le2a/d\le\pi/4$, and the real part of the integrand in
Eq.~\eqref{eq:bessel-integral} is at least $\cos(2a/d)>0$; hence $c_d(a)>0$.

For any constant matrix $C$, the zero-mean property of
$B_a^{(d)}-B_{\mathbf0}^{(d)}$ gives the orthogonal decomposition
\begin{equation}
 \int_{\T^d}\norm{C-B_a^{(d)}(\mathbf q)}_F^2\,d\mathbf q
 =\int_{\T^d}\norm{B_a^{(d)}(\mathbf q)-B_{\mathbf0}^{(d)}}_F^2\,d\mathbf q
  +\norm{C-B_{\mathbf0}^{(d)}}_F^2.
 \label{eq:constant-projection}
\end{equation}
The constant $B_{\mathbf0}^{(d)}=c_d(a)D$ belongs to the allowed family in
Eq.~\eqref{eq:flat-operator-loss}: choose $s=0$, $\lambda=c_d(a)$, and $\phi=0$.
It is therefore the best constant restricted alignment. Equivalently, the normalized
maximum correlation with one global rotation is $c_d(a)$.

Parseval's identity applied entrywise yields
\begin{equation}
 \int_{\T^d}\norm{B_a^{(d)}-B_{\mathbf0}^{(d)}}_F^2
 =\sum_{\mathbf R\ne\mathbf0}\norm{B_{\mathbf R}^{(d)}}_F^2.
\end{equation}
Pointwise unitary invariance gives
$\norm{B_a^{(d)}(\mathbf q)}_F=\norm D_F$. A second use of Parseval, together with
Eq.~\eqref{eq:onsite-average}, therefore gives
\begin{equation}
 \sum_{\mathbf R\ne\mathbf0}\norm{B_{\mathbf R}^{(d)}}_F^2
 =\norm D_F^2-\norm{B_{\mathbf0}^{(d)}}_F^2
 =\norm D_F^2\left[1-c_d(a)^2\right].
\end{equation}
Dividing by $2E_{\mathrm{ref}}^2$ proves Eqs.~\eqref{eq:flat-residual} and
\eqref{eq:flat-tail}. Finally, fixing $s=0$ and $\lambda=1$ leaves the non-onsite
weight plus the onsite mismatch:
\begin{align}
 \inf_\phi L_{\mathrm{op},\phi}^2(0,1)
 &=\kappa[1-c_d(a)^2]+\kappa[1-c_d(a)]^2\\
 &=2\kappa[1-c_d(a)],
\end{align}
which proves Eq.~\eqref{eq:flat-exact-spectrum}.
\end{proof}

\subsection{Orthogonal shell coefficients and factorial tail}
\label{app:proof-orthogonal-tail}

\begin{proof}[Proof of Theorem~\ref{thm:orthogonal-fourier-tail}]
The Jacobi--Anger identity gives
\begin{equation}
 e^{2i\theta_a^{(d)}(\mathbf q)}
 =\prod_{\alpha=1}^d
  \sum_{n\in\mathbb Z}J_n(z)e^{2\pi inq_\alpha}
 =\sum_{\mathbf R\in\mathbb Z^d}
  j_{\mathbf R}(z)e^{2\pi i\mathbf q\cdot\mathbf R}.
\end{equation}
Because $J_{-n}(z)=(-1)^nJ_n(z)$, the cosine coefficient vanishes for odd
$p(\mathbf R)$ and equals $j_{\mathbf R}(z)$ for even $p(\mathbf R)$; the sine
coefficient vanishes for even parity and equals $-ij_{\mathbf R}(z)$ for odd parity.
Substitution into Eq.~\eqref{eq:flat-rotated-pauli} proves
Eq.~\eqref{eq:orthogonal-fourier-block}. Since both Pauli matrices have squared
Frobenius norm two,
\begin{equation}
 \norm{B_{\mathbf R}^{(d)}}_F^2
 =\norm D_F^2\prod_{\alpha=1}^dJ_{R_\alpha}(z)^2.
 \label{eq:block-power-product}
\end{equation}
Equation~\eqref{eq:shell-power} follows by grouping a complete shell. Parseval applied
to the unit-modulus function $e^{iz\sin(2\pi q)}$ gives
$\sum_{n\in\mathbb Z}J_n(z)^2=1$; taking the $d$-fold product proves
Eq.~\eqref{eq:shell-partition}.

If $\norm{\mathbf R}_2^2>S$, at least one coordinate obeys $|R_\alpha|\ge L$.
The union bound and the one-dimensional Bessel partition therefore give
\begin{equation}
 \sum_{\norm{\mathbf R}_2^2>S}
 \prod_{\alpha=1}^dJ_{R_\alpha}(z)^2
 \le 2d\sum_{n=L}^{\infty}|J_n(z)|^2.
 \label{eq:tail-union}
\end{equation}
The defining Bessel series and $(n+m)!\ge n!m!$ imply, for $n\ge0$,
\begin{equation}
 |J_n(z)|
 \le \frac{x^n}{n!}\sum_{m=0}^{\infty}\frac{x^{2m}}{m!}
 =\frac{x^n}{n!}e^{x^2}.
\end{equation}
Moreover, $(L+m)!\ge L!(L+1)^m$, so when $x<L+1$,
\begin{equation}
 \sum_{n=L}^{\infty}\frac{x^{2n}}{(n!)^2}
 \le \frac{x^{2L}}{(L!)^2}
 \sum_{m=0}^{\infty}\left(\frac{x^2}{(L+1)^2}\right)^m
 =\frac{x^{2L}}
 {(L!)^2\left[1-x^2/(L+1)^2\right]}.
\end{equation}
Combining this estimate with Eq.~\eqref{eq:tail-union}, using
$e^{2x^2}=e^{z^2/2}$ and Eq.~\eqref{eq:block-power-product}, proves
Eq.~\eqref{eq:orthogonal-factorial-tail}.
\end{proof}

\subsection{Monotonicity and small-amplitude dimension law}
\label{app:proof-dimension-law}

\begin{proof}[Proof of Corollary~\ref{cor:dimension-law}]
For $x>0$, $J_0'(x)=-J_1(x)$. On $0<x\le\pi/4$, the alternating series
\begin{equation}
 J_1(x)=\sum_{m=0}^{\infty}
 \frac{(-1)^m}{m!(m+1)!}\left(\frac{x}{2}\right)^{2m+1}
\end{equation}
has strictly decreasing term magnitudes because their successive ratio is at most
$\pi^2/128<1$. Hence $J_1(x)>0$. Together with the positivity already proved for
$J_0(2a/d)$, this shows
\begin{equation}
 c_d'(a)=-2J_0(2a/d)^{d-1}J_1(2a/d)<0
\end{equation}
for $a>0$. The threshold monotonicity follows immediately. Finally,
$J_0(x)=1-x^2/4+x^4/64+O(x^6)$. Substituting $x=2a/d$, raising to the $d$th power, and
expanding the two threshold formulas proves
Eqs.~\eqref{eq:correlation-small-a}--\eqref{eq:exact-threshold-small-a}.
\end{proof}

\subsection{Bond-orbit flat-parent comparator}
\label{app:proof-bond-orbit}

\begin{proof}[Proof of Theorem~\ref{thm:bond-orbit}]
Equation~\eqref{eq:bond-orbit-phase} is odd under
$\mathbf q\mapsto-\mathbf q$, so the imaginary part of
Eq.~\eqref{eq:bond-orbit-correlation} integrates to zero. Since
$|2\theta_{a,\mathcal B}|\le\pi/2$, its cosine is nonnegative; it is positive in a
neighborhood of $\mathbf q=\mathbf0$, where the phase vanishes. Hence the integral is
strictly positive. Diagonalizing $D$ by one constant real rotation and repeating
Eq.~\eqref{eq:flat-rotated-pauli} gives
$D_{\mathbf0}^{\mathcal B}=c_{\mathcal B,\boldsymbol\eta}D$. The constant-projection
identity in Eq.~\eqref{eq:constant-projection} and entrywise Parseval then prove the two
residual formulas exactly as in Theorem~\ref{thm:flat-parent}.

For the explicit comparator, apply the Jacobi--Anger expansion to every bond harmonic:
\begin{equation}
 e^{i(2a\eta_j/W)\sin(2\pi\mathbf q\cdot\mathbf b_j)}
 =\sum_{n_j\in\Z}
 J_{n_j}\!\left(\frac{2a\eta_j}{W}\right)
 e^{2\pi i n_j\mathbf q\cdot\mathbf b_j}.
\end{equation}
The finite product of these absolutely convergent series may be integrated termwise.
Torus orthogonality retains exactly the tuples satisfying
$\sum_jn_j\mathbf b_j=\mathbf0$, which proves
Eq.~\eqref{eq:bond-orbit-bessel-sum}. If the integer kernel of the bond matrix is
trivial, only $\mathbf n=\mathbf0$ remains, giving
Eq.~\eqref{eq:independent-bond-product}.
\end{proof}

\subsection{Constrained-series complete-shell certificate}
\label{app:proof-constrained-tail}

\begin{proof}[Proof of Theorem~\ref{thm:constrained-tail}]
Jacobi--Anger expansion and torus orthogonality give
Eq.~\eqref{eq:constrained-coefficient}. The symmetry
$e^{2i\theta(-\mathbf q)}=\overline{e^{2i\theta(\mathbf q)}}$ makes every
$g_{\mathbf R}$ real. Splitting the cosine and sine coefficients into
$(g_{\mathbf R}\pm g_{-\mathbf R})/2$ proves
Eq.~\eqref{eq:constrained-matrix-block}. Hence
\begin{equation}
 \frac{\norm{D_{\mathbf R}^{\mathcal B}}_F^2}{\norm D_F^2}
 =\frac{g_{\mathbf R}^2+g_{-\mathbf R}^2}{2}.
\end{equation}
Parseval for the unit-modulus scalar phase gives
$\sum_{\mathbf R}g_{\mathbf R}^2=1$. Because $\mathcal S_S$ is inversion complete,
the exact matrix tail normalized by $\norm D_F^2$ is
$1-\sum_{\mathbf R\in\mathcal S_S}g_{\mathbf R}^2$. The reverse triangle inequality
$|g_{\mathbf R}|\ge\max\{0,|\widetilde g_{\mathbf R}|-\epsilon_{\mathbf R}\}$
proves Eq.~\eqref{eq:constrained-tail-certificate}.

Solving the triangular constraint gives
$(n_1,n_2,n_3)=(r_1-\ell,r_2+\ell,\ell)$, which proves
Eq.~\eqref{eq:triangular-coefficient}. Taking
$(u,v,w)=(n_4,n_5,n_6)$ in the FCC constraint gives
\begin{equation}
 (n_1,n_2,n_3)=(r_1+u+v,r_2-u+w,r_3-v-w),
\end{equation}
which proves Eq.~\eqref{eq:fcc-coefficient}. Finally, the defining Bessel series gives
$|J_n(z)|\le e^{x^2}x^{|n|}/|n|!$. Summing it over all integers gives $A_0$; summing
over $|n|>K$ and bounding the remaining exponential-series tail geometrically gives
$A_K$. In the triangular remainder the other two Bessel factors have magnitude at most
one. In the FCC remainder, a union bound over the three free indices gives
$3A_KA_0^2$.
\end{proof}


\section{Diagnostic definitions}
\label{app:diagnostics}

Every discrepancy below is defined only after its two arguments have been represented on
the same sampled reciprocal mesh, in the same retained state space, basis, gauge, energy
reference, units, and lattice geometry. A small numerical value cannot repair a missing
state-space identification.

\subsection{Pointwise matrix and sorted-band errors}

Let $A_\ell,B_\ell\in\mathbb C^{r\times r}$ be Hermitian fields evaluated at the same
$M$ reciprocal points, with ordered eigenvalues
$\varepsilon_j^A(\mathbf q_\ell)$ and
$\varepsilon_j^B(\mathbf q_\ell)$. Define the rank-normalized pointwise Hamiltonian
defect
\begin{equation}
 d_{H,\ell}(A,B)=\frac{\norm{A_\ell-B_\ell}_F}{\sqrt r}.
 \label{eq:app-pointwise-hamiltonian-defect}
\end{equation}
The four primary error channels are
\begin{align}
 \epsilon_{H,\max}(A,B)&=\max_\ell d_{H,\ell}(A,B),\\
 \epsilon_{H,\mathrm{rms}}(A,B)
 &=\left[\frac1M\sum_{\ell=1}^M d_{H,\ell}(A,B)^2\right]^{1/2},\\
 \epsilon_{b,\max}(A,B)
 &=\max_{\ell,j}\left|\varepsilon_j^A(\mathbf q_\ell)
                   -\varepsilon_j^B(\mathbf q_\ell)\right|,\\
 \epsilon_{b,\mathrm{rms}}(A,B)
 &=\left[\frac1{Mr}\sum_{\ell=1}^M\sum_{j=1}^r
  \left|\varepsilon_j^A(\mathbf q_\ell)
       -\varepsilon_j^B(\mathbf q_\ell)\right|^2\right]^{1/2}.
 \label{eq:app-primary-error-channels}
\end{align}
They have energy units. A displayed dimensionless value divides them by the declared
$E_{\mathrm{ref}}$. The present retained rank is $r=2$, so the implementation uses the
factor $1/\sqrt2$ in Eq.~\eqref{eq:app-pointwise-hamiltonian-defect}. Pointwise Weyl
control is checked separately; in this normalization,
\begin{equation}
 \epsilon_{b,\max}(A,B)\le
 \sqrt r\,\epsilon_{H,\max}(A,B).
 \label{eq:app-weyl-check}
\end{equation}
No equality between matrix and band errors is assumed.

\subsection{Interpolation, deliberate truncation, and total error}

Write $H^{\mathrm{dir}}$ for the direct withheld-mesh field,
$H_N^{\mathrm{ord}}$ for the full ordinary Fourier interpolant, and
$H_{N,R}^{\mathrm{ord}}$ for its complete-radius truncation. For each of the four
functionals in Eq.~\eqref{eq:app-primary-error-channels}, the three approximation
channels are
\begin{align}
 \mathcal E_{\mathrm{interp}}&=\mathcal E(H_N^{\mathrm{ord}},H^{\mathrm{dir}}),\\
 \mathcal E_{\mathrm{trunc}}&=\mathcal E(H_{N,R}^{\mathrm{ord}},H_N^{\mathrm{ord}}),\\
 \mathcal E_{\mathrm{total}}&=\mathcal E(H_{N,R}^{\mathrm{ord}},H^{\mathrm{dir}}).
 \label{eq:three-error-channels}
\end{align}
These channels are not added: interpolation and truncation may reinforce or cancel. For
a strain-responsive control, the same definitions are applied to the gauge-aligned
difference field $\Delta H_s$ after constructing it in one common retained space. The
geometry-only strain control reports locality under a changing metric and is not a
three-channel strain-potential calculation.

\subsection{Fourier power, complete shells, and first passing}

For ordinary hopping blocks $T_{\mathbf n}$ and a complete physical-metric support
$\mathcal S_s$, define
\begin{align}
 P_{\mathrm{shell}}(s)
 &=\frac1r\sum_{\mathbf n\in\mathcal S_s\setminus\mathcal S_{s-1}}
   \norm{T_{\mathbf n}}_F^2,\\
 P_{\mathrm{omit}}(s)
 &=\frac1r\sum_{\mathbf n\notin\mathcal S_s}
   \norm{T_{\mathbf n}}_F^2.
 \label{eq:app-shell-power}
\end{align}
Division by $E_{\mathrm{ref}}^2$ gives the dimensionless power. Cyclic Nyquist power is
reported separately and is not silently reassigned to an ordinary conjugate-closed
hopping block. Parseval defects compare reciprocal-space power with the complete cyclic
block inventory; truncation Parseval defects compare deliberate omitted power with its
reciprocal-space reconstruction consequence.

For diagnostic channel $X$ and tolerance $\tau$, the sampled first-passing shell is
\begin{equation}
 \widehat s_X(\tau)=
 \min\{s:\epsilon_{X,\mathrm{total},\max}(s)\le\tau\}.
 \label{eq:app-first-passing}
\end{equation}
The record retains the physical radius, block count, passing error, and immediately
preceding error. If no complete shell passes before the declared cap, the decision is
unresolved rather than assigned to the cap.

\subsection{Alignment, reconstruction, and isolation diagnostics}

The reconstruction ledger keeps parent Hermiticity, complete-eigensystem reconstruction,
frame orthonormality, projector reconstruction, parent--retained lifting, Fourier
reconstruction, covariant equality, and restricted global alignment as separate
defects. External retained-subspace isolation is reported both as the minimum sampled
gap and as a continuous lower bound. A sampled gap is not relabeled as a continuous
certificate.

Restricted alignment records one global rotation, the optimized nonnegative splitting
scale, the finite-mesh maximum correlation, the minimum operator threshold, and the
distinct unit-splitting exact-spectrum joint threshold. A fixed-unit-scale rotation fit
cannot substitute for the optimized-scale threshold. Flat-family analytic correlations,
quadrature defects, residual--tail identities, amplitude midpoint defects, and
amplitude-resolution ladders remain separate checks.

\subsection{Strain candidates, sampled envelopes, and transition brackets}

At sampled strain $s_i$, let $\kappa_i$ and $\mu_i$ denote the restricted-alignment
splitting scale and maximum correlation. The pointwise optimized requirement and the
fixed reference-candidate requirement are
\begin{align}
 \tau_{o,\min}(s_i)
 &=\sqrt{\kappa_i[1-\max(0,\mu_i)^2]},\\
 \tau_{o,p_0}(s_i)
 &=\sqrt{2\kappa_i(1-\mu_i)},
 \qquad p_0=(0,1).
 \label{eq:app-strain-candidate-requirements}
\end{align}
The first quantity permits an independently optimized candidate at every sampled
strain. It is not one common witness. The second evaluates one prescribed parameter
path, but it is not a witness until declared spectral and operator limits establish
feasibility. Their hatted maxima, and the hatted locality radius and block-count maxima,
are finite-sample statistics:
\begin{equation}
 \widehat\tau_{o,\min}=\max_i\tau_{o,\min}(s_i),\quad
 \widehat\tau_{o,p_0}=\max_i\tau_{o,p_0}(s_i),\quad
 \widehat R_\varepsilon=\max_iR_\varepsilon(s_i),\quad
 \widehat N_\varepsilon=\max_iN_\varepsilon(s_i).
 \label{eq:app-strain-sampled-envelopes}
\end{equation}
Each maximum retains its own sample location; the locations need not agree.

The response-sweep qualification envelope additionally retains
\begin{equation}
 \min_i g_{\mathrm{ext}}(s_i),\qquad
 \min_i\sigma_{\mathrm{seed}}(s_i),\qquad
 \min_i\sigma_{\mathrm{frame}}(s_i),\qquad
 \max_i r_{\mathrm{point}}(s_i),\qquad
 \max_i r_{\mathrm{global}}(s_i),
 \label{eq:app-strain-qualification-envelope}
\end{equation}
where $g_{\mathrm{ext}}$ is the sampled external isolation gap,
$\sigma_{\mathrm{seed}}$ is the minimum seed-projector overlap,
$\sigma_{\mathrm{frame}}$ is the minimum singular value of the pointwise frame overlap,
and the two $r$ channels are pointwise and restricted-global frame-alignment residuals.
Every extremum retains the first sampled strain at which it occurs. These quantities
qualify the retained state-space identification; they are not response observables.

For the frozen-parent and radial-response controls, the sweep also takes separate sampled
maxima of all four errors in Eq.~\eqref{eq:app-primary-error-channels} for the aligned
response, interpolation, deliberate truncation, and total channels. The same reduction
is applied to $\Delta H_s$, ordinary omitted power, cyclic Nyquist power, scalar and
traceless response components, onsite and offsite response power, and minimum and
maximum prescribed hopping-response factors. No maxima from distinct channels are
combined into a synthetic sample, and none is relabeled as a continuum supremum.

A detected transition bracket is an adjacent sampled interval whose endpoint values of
a discrete channel differ. The present discrete channels are first-passing block count
at each tail tolerance and number of distinct bond-length orbits. Five bisections reduce
a base-grid interval of width at most $1/200$ to width at most $1/6400$. The bracket
localizes a detected change only. Equal endpoint values do not prove that the channel is
constant inside the interval, and the bracket supplies no derivative or interval bound
for the continuous path.

\subsection{Transport and closure diagnostics}

For rank-two closure diagnostics, let $W_\alpha(\mathbf k_\perp)$ be the ordered product
of full-rank transport links around reciprocal direction $\alpha$ at transverse base
point $\mathbf k_\perp$. If
$-\pi\le\phi_{\alpha,0}(\mathbf k_\perp)\le
\phi_{\alpha,1}(\mathbf k_\perp)\le\pi$ are its ordered principal eigenphases, define
\begin{equation}
 \Phi_j=\max_{\alpha,\mathbf k_\perp}
 |\phi_{\alpha,j}(\mathbf k_\perp)|,
 \qquad j\in\{0,1\}.
 \label{eq:multidimensional-closure-phase-reduction}
\end{equation}
This reduction is applicable only when every contributing transport link is full rank.
It is accompanied by the minimum transport singular value, Procrustes residual,
rank-cutoff and roundoff controls, link-unitarity defect, closure-unitarity defect, and
maximum plaquette phase. These closure and plaquette channels have no asserted
topological meaning here. A rank-deficient link makes the dependent closure channels
explicitly inapplicable rather than numerically zero.

\subsection{Numerical validation boundary}

The numerical checks also retain mesh separation, metric covariance, basis-distance
extrema where applicable, refinement changes, and independent reconstruction defects.
Agreement of these diagnostics verifies the declared finite calculation. It does not
establish a continuum limit, supply an externally chosen adequacy tolerance, or convert
a synthetic model into a material-specific Hamiltonian.


\section{Coordinate, Fourier, and state-space conventions}
\label{app:conventions}

Let the columns of $A\in\mathbb R^{d\times d}$ be direct primitive vectors. Reduced
reciprocal coordinates $\mathbf q\in\T^d$ and integer lattice modes
$\mathbf n\in\Z^d$ are mapped to Cartesian vectors by
\begin{equation}
 B=2\pi A^{-\mathsf T},\qquad
 \mathbf k=B\mathbf q,\qquad
 \mathbf R=A\mathbf n.
 \label{eq:app-direct-reciprocal-map}
\end{equation}
Consequently $\mathbf k\cdot\mathbf R=2\pi\mathbf q\cdot\mathbf n$. The
nonorthogonal protocols retain this reduced-coordinate convention; the physical radius
of a hopping mode is $\norm{A\mathbf n}$ rather than a coordinate-axis norm.

For a matrix field $H(\mathbf q)$, the continuous Fourier convention is
\begin{equation}
 H_{\mathbf n}=\int_{\T^d}
 e^{-2\pi i\mathbf q\cdot\mathbf n}H(\mathbf q)\,d\mathbf q,
 \qquad
 H(\mathbf q)=\sum_{\mathbf n\in\Z^d}
 e^{2\pi i\mathbf q\cdot\mathbf n}H_{\mathbf n}.
 \label{eq:app-fourier-convention}
\end{equation}
A finite even mesh first produces a cyclic centered inventory. Cyclic Nyquist modes are
reported separately before selecting an ordinary conjugate-closed hopping inventory.
Every finite support is completed in the physical metric and retains all equal-radius
ties.

Hamiltonian subtraction is performed only after state spaces have been identified. The
parent orbital field, complete eigensystem, retained projector, periodic retained frame,
and Wannier blocks remain distinct representations. Pointwise equality of sorted
eigenvalues is not used as a substitute for an operator identification.

\section{Model inventory and numerical design}
\label{app:protocols}

The fixed-geometry comparison contains four model families, and the strain study adds
two geometry paths.

\begin{table}[htbp]
\centering
\small
\caption{Model and control inventory.}
\label{tab:app-protocol-inventory}
\begin{tabular}{@{}>{\raggedright\arraybackslash}p{0.30\textwidth}
                    >{\raggedright\arraybackslash}p{0.62\textwidth}@{}}
\toprule
Family or path & Geometry and role \\
\midrule
Square 2-D & Orthogonal directed-parent baseline \\
Triangular/honeycomb 2-D & Nonorthogonal target with matched square and triangular
Bravais controls \\
Simple-cubic 3-D & Orthogonal directed-parent baseline \\
FCC/diamond 3-D & Nonorthogonal target with matched simple-cubic and FCC Bravais
controls \\
Rhombic strain & Unit-area triangular-to-square path with three controls \\
Bain strain & Unit-volume FCC-to-BCC path with three controls \\
\bottomrule
\end{tabular}
\end{table}

The construction and shifted evaluation meshes are tensor products. The evaluation
shift is one third of one evaluation-grid interval in every reduced reciprocal
coordinate. Table~\ref{tab:app-numerical-design} gives the three-level mesh ladder. The
reported $10^{-3}$ and $10^{-5}$ tail-error tolerances select first-passing complete
supports; they are support-selection coordinates, not acceptance limits for restricted
operator mismatch.

\begin{table}[htbp]
\centering
\small
\caption{Fixed-geometry numerical design. Mesh sizes are per reciprocal coordinate.}
\label{tab:app-numerical-design}
\begin{tabular}{@{}lccc@{}}
\toprule
Family & Construction sizes & Evaluation sizes & Shell cap \\
\midrule
Square 2-D & $16,32,64$ & $33,65,129$ & 25 \\
Triangular/honeycomb 2-D & $16,32,64$ & $33,65,129$ & 21 \\
Simple-cubic 3-D & $8,16,32$ & $17,33,65$ & 9 \\
FCC/diamond 3-D & $8,16,32$ & $17,33,65$ & 7 \\
\bottomrule
\end{tabular}
\end{table}

The orthogonal gauge amplitude is sampled at $0$, $\pi/8$, and $\pi/4$. The
triangular/honeycomb family uses maximum amplitude $\pi/(4\sqrt6)$ and the FCC/diamond
family uses $\pi/(4\sqrt{12})$; both include zero and the midpoint. Every family samples
parent scales $\rho\in\{0,1\}$. These are synthetic controls, not graphene or silicon
parameterizations.

\section{Cross-protocol requirement coordinates}
\label{app:tolerance-requirements}

The positive-splitting-scale optimization and the exact-spectrum unit-scale question are
kept separate. The earlier two-coordinate requirement reduction records
\begin{equation}
 x=\frac{\tau_{o,\min}}{\sqrt\kappa},\qquad
 y=\frac{\tau_{o,\mathrm{joint}}}{E_{\mathrm{ref}}}.
 \label{eq:app-requirement-pair}
\end{equation}
The componentwise maxima over cases, parent scales, and amplitudes need not originate
from the same sampled point and are not attainable decisions in general.
Table~\ref{tab:app-case-requirements} retains these scalar audit coordinates for each
case, but the main text replaces their visual envelope with the shell--scale Pareto
formulation in Eq.~\eqref{eq:pareto-decision-objectives}.

\begin{table}[htbp]
\centering
\small
\caption{Componentwise case envelopes for the sampled fixed-geometry calculation.
These coordinates summarize calculated requirements; they are not acceptance limits.}
\label{tab:app-case-requirements}
\begin{tabular}{@{}lcc@{}}
\toprule
Case & $\max\tau_{o,\min}/\sqrt\kappa$ &
$\max\tau_{o,\mathrm{joint}}/E_{\mathrm{ref}}$ \\
\midrule
Square directed parent & 0.688457 & 0.328854 \\
Matched square Bravais control & 0.586917 & 0.273379 \\
Triangular Bravais control & 0.584676 & 0.263702 \\
Honeycomb-inspired target & 0.468419 & 0.083542 \\
Simple-cubic directed parent & 0.584781 & 0.270259 \\
Matched simple-cubic Bravais control & 0.432793 & 0.195029 \\
FCC Bravais control & 0.430770 & 0.189573 \\
Diamond-inspired target & 0.347324 & 0.052890 \\
\bottomrule
\end{tabular}
\end{table}

Table~\ref{tab:app-case-requirements} reports requirement coordinates without deriving
acceptance limits from the observed values. Operating-point selection requires an
external application criterion.

\section{Case-resolved shell--scale decision landscapes}
\label{app:case-decision-landscapes}

Table~\ref{tab:app-decision-landscape-inventory} records the fixed-geometry endpoint
sampling used in the case-resolved plots. Each case uses its maximum declared parent
scale and gauge amplitude, every complete shell under its own physical metric and shell
cap, and 17 uniformly spaced values of the scale-path coordinate in
Eq.~\eqref{eq:scale-path}. Exact componentwise dominance is evaluated only among these
calculated vertices. The number of nondominated vertices is therefore a sampling
diagnostic, not a continuum invariant or a measure of scientific adequacy.

\begin{table}[htbp]
\centering
\small
\caption{Sampled shell--scale landscape inventory.}
\label{tab:app-decision-landscape-inventory}
\begin{tabular}{@{}lrrr@{}}
\toprule
Case & $\lambda_*$ & Calculated & Nondominated \\
\midrule
Square directed parent & 0.729890 & 238 & 120 \\
Matched square Bravais control & 0.813169 & 221 & 127 \\
Triangular Bravais control & 0.835297 & 187 & 156 \\
Graphene-inspired honeycomb & 0.883506 & 187 & 126 \\
Simple-cubic directed parent & 0.812597 & 153 & 85 \\
Matched simple-cubic Bravais control & 0.902347 & 119 & 80 \\
FCC Bravais control & 0.913939 & 136 & 116 \\
Silicon-inspired diamond & 0.937745 & 136 & 34 \\
\bottomrule
\end{tabular}
\end{table}

The parent scale is $1$ in every row. The gauge amplitude and complete-shell inventory
are as follows; parenthesized block counts give $N(S)$ in the same order as the listed
shell indices.
\begin{description}
 \item[Square directed parent:] $a=\pi/4$;
 $S=(0,\allowbreak 1,\allowbreak 2,\allowbreak 4,\allowbreak 5,\allowbreak 8,\allowbreak 9,\allowbreak 10,\allowbreak 13,\allowbreak 16,\allowbreak 17,\allowbreak 18,\allowbreak 20,\allowbreak 25)$;
 $N=(1,\allowbreak 5,\allowbreak 9,\allowbreak 13,\allowbreak 21,\allowbreak 25,\allowbreak 29,\allowbreak 37,\allowbreak 45,\allowbreak 49,\allowbreak 57,\allowbreak 61,\allowbreak 69,\allowbreak 81)$.
 \item[Matched square control:] $a=\pi/(4\sqrt6)$;
 $S=(0,\allowbreak 1,\allowbreak 2,\allowbreak 4,\allowbreak 5,\allowbreak 8,\allowbreak 9,\allowbreak 10,\allowbreak 13,\allowbreak 16,\allowbreak 17,\allowbreak 18,\allowbreak 20)$;
 $N=(1,\allowbreak 5,\allowbreak 9,\allowbreak 13,\allowbreak 21,\allowbreak 25,\allowbreak 29,\allowbreak 37,\allowbreak 45,\allowbreak 49,\allowbreak 57,\allowbreak 61,\allowbreak 69)$.
 \item[Triangular control:] $a=\pi/(4\sqrt6)$;
 $S=(0,\allowbreak 1,\allowbreak 3,\allowbreak 4,\allowbreak 7,\allowbreak 9,\allowbreak 12,\allowbreak 13,\allowbreak 16,\allowbreak 19,\allowbreak 21)$;
 $N=(1,\allowbreak 7,\allowbreak 13,\allowbreak 19,\allowbreak 31,\allowbreak 37,\allowbreak 43,\allowbreak 55,\allowbreak 61,\allowbreak 73,\allowbreak 85)$.
 \item[Honeycomb:] $a=\pi/(4\sqrt6)$; it uses the same shell indices and block counts
 as the triangular control.
 \item[Cubic parent:] $a=\pi/4$;
 $S=(0,1,2,3,4,5,6,8,9)$; $N=(1,7,19,27,33,57,81,93,123)$.
 \item[Matched cubic control:] $a=\pi/(4\sqrt{12})$;
 $S=(0,1,2,3,4,5,6)$; $N=(1,7,19,27,33,57,81)$.
 \item[FCC control:] $a=\pi/(4\sqrt{12})$;
 $S=(0,1,2,3,4,5,6,7)$; $N=(1,13,19,43,55,79,87,135)$.
 \item[Diamond:] $a=\pi/(4\sqrt{12})$; it uses the same shell indices and block
 counts as the FCC control.
\end{description}

The axes of every plot are the three components of
$\mathcal J_s(S,u)$ in Eq.~\eqref{eq:pareto-decision-objectives}, evaluated here at the
fixed-geometry endpoint amplitude rather than along a strain path. Color gives the
scale-path coordinate $u$. The displayed triangular faces connect only adjacent cells in
the sampled $(S,u)$ grid. They do not assert that the objective-space image is convex,
smooth, continuously covered, or valid between vertices.

\begin{figure}[p]
 \centering
 \includegraphics[width=0.92\textwidth]{figures/square-directed-parent-shell-scale-decision-landscape.pdf}
 \caption{Square directed-parent decision landscape. The comparatively broad separation
 between $\lambda_*$ and unity produces a visible operator--spectrum scale compromise.}
 \label{fig:app-decision-square-parent}
\end{figure}

\begin{figure}[p]
 \centering
 \includegraphics[width=0.92\textwidth]{figures/matched-square-bravais-control-shell-scale-decision-landscape.pdf}
 \caption{Matched square Bravais-control decision landscape. This control changes the
 parent law while retaining the square metric.}
 \label{fig:app-decision-square-control}
\end{figure}

\begin{figure}[p]
 \centering
 \includegraphics[width=0.92\textwidth]{figures/triangular-bravais-control-shell-scale-decision-landscape.pdf}
 \caption{Triangular Bravais-control decision landscape. Its shell inventory is defined
 by the non-orthogonal physical metric rather than square-coordinate radius.}
 \label{fig:app-decision-triangular-control}
\end{figure}

\begin{figure}[p]
 \centering
 \includegraphics[width=0.92\textwidth]{figures/graphene-inspired-honeycomb-shell-scale-decision-landscape.pdf}
 \caption{Graphene-inspired honeycomb decision landscape. ``Inspired'' denotes the
 synthetic geometry and connectivity, not a graphene material parametrization.}
 \label{fig:app-decision-honeycomb}
\end{figure}

\begin{figure}[p]
 \centering
 \includegraphics[width=0.92\textwidth]{figures/simple-cubic-directed-parent-shell-scale-decision-landscape.pdf}
 \caption{Simple-cubic directed-parent decision landscape.}
 \label{fig:app-decision-cubic-parent}
\end{figure}

\begin{figure}[p]
 \centering
 \includegraphics[width=0.92\textwidth]{figures/matched-simple-cubic-bravais-control-shell-scale-decision-landscape.pdf}
 \caption{Matched simple-cubic Bravais-control decision landscape.}
 \label{fig:app-decision-cubic-control}
\end{figure}

\begin{figure}[p]
 \centering
 \includegraphics[width=0.92\textwidth]{figures/fcc-bravais-control-shell-scale-decision-landscape.pdf}
 \caption{FCC Bravais-control decision landscape using complete FCC metric shells.}
 \label{fig:app-decision-fcc-control}
\end{figure}

\begin{figure}[p]
 \centering
 \includegraphics[width=0.92\textwidth]{figures/silicon-inspired-diamond-shell-scale-decision-landscape.pdf}
 \caption{Silicon-inspired diamond decision landscape. Its optimized scale is closest to
 unity and its sampled nondominated set is the most selective of the eight cases. This
 remains a synthetic rank-two control, not a silicon material calculation.}
 \label{fig:app-decision-diamond}
\end{figure}

\clearpage
\section{Finite-strain control definitions}
\label{app:strain-controls}

The strain sequence separates three distinct questions. It does not assign the geometry
control, null response, and prescribed response to a single model.

\subsection{Geometry-only control}

The 2-D unit-area rhombic path is
\begin{equation}
 A_2(\theta)=\frac{1}{\sqrt{\sin\theta}}
 \begin{pmatrix}1&\cos\theta\\0&\sin\theta\end{pmatrix},
 \qquad \theta:\pi/3\longrightarrow\pi/2.
 \label{eq:app-rhombic-path}
\end{equation}
The 3-D path is the unit-primitive-volume Bain deformation from FCC to BCC in a common
body-centered tetragonal primitive basis. In both cases, the reduced-coordinate parent,
imposed gauge, and Fourier-block weights remain fixed. The changing physical metric
alone reorders complete shells and changes first-passing physical radius and retained
block count.

\subsection{Frozen-parent null-response control}

The frozen-parent control uses the four-state bipartite parent and the
same imposed gauge at every strain point while keeping onsite and hopping matrices fixed.
Its gauge-aligned strain difference is expected to vanish up to numerical roundoff. This
is a null-response control, not a strain-dependent Hamiltonian model.

\subsection{Synthetic radial-response control}

The radial-response control changes only the selected inter-sublattice bond amplitudes:
\begin{equation}
 \frac{t_j(s)}{t_j(0)}=
 \exp\!\left[-\left(\frac{d_j(s)}{d_j(0)}-1\right)\right].
 \label{eq:app-radial-response}
\end{equation}
No onsite response, refit, structural relaxation, or material calibration is included.
The response field is aligned in a common retained space before scalar/traceless,
onsite/offsite, interpolation, deliberate-truncation, and total channels are reported.
The interpolation, truncation, and total channels are not added.

At the final 2-D metric, the calculated response factors span
$0.81758$--$1.27146$ and the maximum aligned retained-space perturbation is
$0.063575E_{\mathrm{ref}}$. At the BCC endpoint, the four selected bonds share response
$0.97564$ and the maximum perturbation is $0.007289E_{\mathrm{ref}}$. These are
results for the prescribed synthetic law, not material parameters or predictions.

\section{Independent numerical reconstruction and scope}
\label{app:verification}

An independent numerical oracle reconstructs mesh coordinates, minimum
construction--evaluation separation, parent and retained fields, restricted global
alignment, direct finite-sum Fourier blocks, Parseval identities, Hermiticity defects,
complete physical shells, interpolation/truncation/total matrix and band channels, Weyl
bounds, and first-passing shell decisions. It does not invoke the production Fourier
implementation it verifies. Nonorthogonal cases retain reduced coordinates and integer
modes while physical radii use $A\mathbf n$.

Agreement between the two implementations checks numerical reconstruction but does not
turn a finite mesh into a continuum result or a synthetic parent into a material model.
The geometry-only, frozen-parent, and radial-response calculations are therefore
interpreted within their declared finite meshes, shell caps, and response laws.


\bibliographystyle{unsrt}
\bibliography{references}

\end{document}
